% !TEX root = ../cubic-moduli-descent.tex
\section{構造対象付きモジュライと降下の主要結果}
\label{sec:moduli-overview}

本節では、後続の証明に用いる仮定と主要な結論を、
満支持複数ブロック型、局所一ブロック型、比較核の降下の順にまとめる。
各主張の直後に詳しい定理と証明の所在を示す。
全 intrinsic 対象の表示分類や、任意の滑らかな非可換三次の
可換 Kuznetsov 実現を結論するものではない。

\subsection{満支持複数ブロック型で成立すること}
$n=3k-2$, $k\geq2$ とし、$S_i$ は滑らかな点曲線 $E_i$ と
平行移動 $\sigma_i=t_{p_i}$ を持つ三次元楕円 Sklyanin 代数とする。
$p_i,2p_i,3p_i\ne0$、$j(E_i)\ne0,1728$、$j(E_i)$ は相異なると仮定する。
各正則中心三次を $\Omega_i$ とし、
\[
 B=\bigotimes_{i=1}^k S_i,\qquad
 A=B/\left(\sum_i a_i\Omega_i\right),\qquad a_i\ne0
\]
と置く。定理\ref{mbs:thm:presentation-stack}は、共通 grading scalar を除いた
全次数付き表示の滑らかな既約 DM stack $\mathfrak P_k$ を構成する。
その次元は $2k$、幾何点の慣性は
$(\prod_i\mathcal H_{27,i})/\Delta\mu_3$、位数は $3^{3k-1}$ である。
二ブロック三次四様体では、次元四と位数 $243$ を得る。
この結果は相対 Jacobian による各因子の表示と、
次数付き代数からのブロック・中心三次直線の認識を用いる。

一方、点スキームは $\coprod_iE_i$ であり、その双線型イデアルは
実際の二次関係より $9\binom{k}{2}$ 次元大きい
（定理\paperref{main}{tbp:thm:point-graph}、系\paperref{main}{tbp:cor:all-blocks}）。
従って点曲線を知るだけではブロック間の交換関係を選べない。
ブロック別 $\mu_3$ による Zhang twist は、
同じ $(3k-3)$-Veronese と同じ $\OO$-marked 圏を持ちながら、
元の可換 tensor 積表示と次数付き同型でない例を与える。
Zhang の圏同値は既知の入力\cite{TBPZhangTwist}であり、
ここでは満支持型の中心商に適用して次数一復元の具体的な限界を示す。

相対双対では基底直線束
$\mathscr D=\det(\bigoplus_iW_i)\otimes\ell^\vee$ を保持する。
定理\paperref{main}{mbsd:thm:serre}と系\paperref{main}{mbsd:cor:saturation}により、
Serre 関手は $(-3k+3)[3k-2]\otimes\mathscr D$ であり、全整数の切断は飽和する。
全構造対象保存同値が認識する反標準切断環は
$\bigoplus_{q\geq0}A_{(3k-3)q}\otimes\mathscr D^{-q}$ である。
さらに $jp_i\ne0$ $(1\leq j\leq3k-3)$ の有限個の条件の下で、
第三 Veronese の中心を直接計算し、中心平坦性を証明する。
定理\ref{mbv:thm:all-o-third}は fpqc 層として
\[
 \operatorname{Eq}_{\OO}(\operatorname{QGr}A,\operatorname{QGr}A')
 \simeq
 \underline{\operatorname{Isom}}_{\rm gr}(A^{(3)},(A')^{(3)})/\mathbf G_m
\]
を与える。二ブロックでは第三 Veronese は $55$ 生成元と $2619$ 個の独立な
二次関係を持ち、固定模型間の全同値層は有限表示のスキームで表される。
中心一次部分を殺した商 $D=C/(Z(C)_1)$ の点加群は、各因子の第三 Veronese に由来する。
その純成分のイデアル $I_i$ と余法加群 $I_i/I_i^2$ は、元の shift の逆を復元する。
従って、可換 tensor 表示の幾何点では全 $\OO$-同値から
$(E_i,L_i,\sigma_i)$ と次数付き同型類が決まる
（定理\ref{mbiq:thm:triple-recovery}、系\ref{mbiq:cor:commuting-o-rigidity}）。
本稿ではさらに、各 $I_i$ の冪の右飽和から $D$ の全多重次数を回復する。
各剰余次数加群の次数付き自己準同型環が $K$ であることを使うと、
$D=(\bigotimes_iT_i)^{(3)}$ の全自己同型は $\bigotimes_iT_i$ へ持ち上がる
（定理\ref{mbca:thm:center-automorphisms}）。
一方、中心完全交叉の AS--Cohen--Macaulay 性により
$\overline M_j=(E(j))^{(3)}$ の負次数 $\Hom,\Ext^1$ が消える。
これを全中心変数による thickening に適用し、正次数の次数別安定性で
形式的同型を実際の次数付き同型へ戻す
（補題\ref{mbcr:lem:graded-lifting}）。
この二つの議論を接続した定理\ref{mbca:thm:full-inertia}は
\[
 \Aut_{\rm gr}(A^{(3)})/K^\times=\mathcal K_k,
 \qquad
 \Aut_{\OO}(\QGr A)=\mathcal K_k
\]
を与える。後者には上記の有限個の捩れ条件を課し、前者には冒頭の基本条件だけを使う。
中心還元の射影的な核も自明である（定理\ref{mbcr:thm:reduction-rigidity}）。
したがって二ブロックの全慣性の位数は正確に $243$ となる。
さらに、一次中心変形が自明なら中心変数の符号反転が還元核を与えることから、
二ブロックでは $\kappa_C\ne0$ を得る（系\ref{mbcr:cor:nonzero-class}）。

可換三次への支配的特殊化と関係空間の generic flatness による独立の構成から、
非空開上の全 $\OO$-局所型は滑らかな既約 DM stack $\mathfrak M_k^\circ$ となり、
次元は $2k$ である（定理\ref{mbg:thm:generic-stack}）。
さらに、射影枠を使って $\mathfrak P_k^\circ\to\mathfrak M_k^\circ$ の表現可能性を示し、
標数零の一般滑らかさと有限表示降下から、非空開上で有限 \'etale とする。
全慣性の一致によって階数は一となるため、この射は同型である
（系\ref{mbrr:cor:stack-isomorphism}）。
従ってこの開では、非被約底を含む全ての試験基底上で
全 $\OO$-同値の層と次数付き同型のスカラー商が一致する。
全慣性の中心は、順序付きブロック被覆上で
$(\mu_3^k)/\mu_3^{\rm diag}$ と内在的に回復される。
全境界や任意の抽象 Abel 圏変形を覆うことは含めない。

共通スカラーを戻した表示 gerbe の類 $\beta$ は、各非空開で正確に位数三である
（定理\ref{mbsd2:thm:gerbe-period}）。
上界は関係直線 $\ell$ の重み三、下界は実表示慣性の commutator に含まれる
共通 $\mu_3$ による。順序付き被覆では重み一の最小ベクトル束階数も三である。
$\OO(j)$ の重みは $j$ であり、例外的成分は
$\operatorname{Perf}(\mathfrak P_k,-j\beta)$ として降下する。
RHom と対象の重みの相殺により compact 残余 projector 自体は降下する
（定理\ref{mbsd2:thm:projector-descent}）。
標準 rotation は重み一であり、内部 shift を $\rho$ へ移す MF/window 比較がある範囲では、関係直線を含む周期
$\rho^3\otimes\ell^\vee\simeq[2]$ と、方程式を標識した後の周期を区別する
（命題\ref{mbsd2:prop:line-period}）。
これは line-valued MF の既知の翻訳規約
\cite[Definitions~1.1, 1.3]{MBS2PolishchukVaintrob}を本表示へ適用したものである。

中心群 $Z_0$ の作用は、共通 scalar の自然同型を用いて合成則ごと降下する
（定理\paperref{main}{mbza:thm:central-action}）。
一方、その線型的持上げの選択は
$\{(a_i)\in\mathbf F_3^k:\sum_i a_i=1\}$ という
$\widehat Z_0$-torsor であり、基準変更は rotation の指標 twist となる。
ブロック置換に不変な選択は $3\nmid k$ のときに限り存在する
（系\paperref{main}{mbza:cor:permutation-obstruction}）。
この有限群の障害は中心作用自体を妨げず、表示 gerbe の $\beta$ とも別である。
群作用・同変化・双対群作用の一般論は Elagin\cite{MEDElagin}、
随伴と半直交分解の同変性は Shinder\cite{MBZAShinder}を入力とし、
ここでは scalar 自明化と全ての指標変換を明示する。

因子の中心係数を忘れる射は $Z_0$-gerbe であり、
rigidification $\mathfrak P_k\mathbin{/\mkern-6mu/}Z_0$ は
楕円 Sklyanin 因子の組の stack $\mathfrak B_k$ と同型になる
（命題\ref{mbfm:prop:rigidification}）。
これは Romagny\cite[Theorem~5.1]{MBFMRomagny}の一般的な操作を、
係数 torus 上の三乗写像によって具体化したものである。
平面三次の組の stack $\mathfrak E_k$ への射
$\mathfrak B_k\to\mathfrak E_k$ は表現可能・滑らか・全射で相対次元 $k$ となり、
曲線積の普遍族を引き戻せる（命題\ref{mbfm:prop:point-map}、
命題\ref{mbfm:prop:product-family}）。
中心作用付き圏の同変化は、この rigidification 上の圏としても記述できるが、
それだけで曲線積との同値や表示 gerbe の中立化を結論しない。

根については、$\mathsf s=(1)$ としたとき
$\mathsf r\OO\simeq\OO(1)$ を満たす自己同値は一意に
$\mathsf r\simeq\mathsf s\Phi_h$、$h\in\mathcal K_k$ と書け、
$\mathsf r^3\simeq\mathsf s^3$ も自動的に成立する
（系\ref{mbca:cor:marked-roots}）。
その自然同型類は $3^{3k-1}$ 個、全 $\OO$-群による共役類は
$(11^k+2\cdot8^k)/3$ 個であり、二・三ブロックでは $83,785$ となる。
この関手の分類には $6p_i\ne0$ を追加しない。
一つの次数付き自己同型の冪で定まる全代数的 Zhang twist の表示を
直接分類する定理\ref{mbr:thm:algebraic-twist-groupoid}では、この追加条件を使う。
両端が可換 tensor 積表示である一般 twisting system の剛性も保持する
（定理\ref{mbr:thm:commuting-twist-rigidity}）。
根の三乗同型そのものや coherence は標識しておらず、
$\mathsf r^3\simeq\mathsf s^3$ だけから $\mathsf r\OO$ を決定する問題は残る。

Veronese の自己同型持上げについて、Aitzhanova--Makar-Limanov--Umirbaev
\cite[Theorem~1, Corollaries~4, 5]{MBCAAMU}は自由 Poisson 代数と多項式環を扱う。
同論文の $d$-graded は次数を $d$ で剰余化した意味であり、通常の次数保存とは区別する。
また Chan--Young--Zhang\cite[Theorems~0.1, 0.2]{MBCACYZ}は、
パラメータが非自明な1の冪根である skew polynomial 環の Veronese を判別式で扱う。
これらは本稿の楕円 tensor 中心商への直接の持上げ定理ではない。
本稿の認識は、余法加群・飽和冪・負次数 Ext と中心平坦性を組み合わせる。
THCR の理論、Heisenberg 群、次数付き持上げの一般原理自体を新規成果とはせず、
本件への組合せと適用については初出の網羅的確認を継続する。

多ブロック一変数族では、既存の定理\paperref{main}{cft:thm:complete}と
定理\paperref{main}{gdr:thm:generic}を入力として、
$G=(\mu_3)^{k-1}$ 不変な低 rotation 積の積分空間から
分裂した復元方程式 $F_1+\cdots+F_k$ を元の関数体上に得る。
同変核は crossed product 上の perfect 加群として coherence ごと有限型底へ降ろし、
その閉点において有限拡大上の $G$ 同変・rotation・指定周期保存比較を得る
（定理\paperref{main}{med:thm:multiblock-generic}）。
これは指定した形式核の全 jet の代数化ではない。
可換側では Ballard--Favero--Katzarkov の Morita 積公式
\cite[Definition~3.50, Corollary~5.18]{NFBFKKernels}を
$\Gamma=\{(\lambda_i):\lambda_i^3\text{ が共通}\}$ に適用する。
その Orlov 比較\cite[Theorem~6.13]{NFBFKKernels}と合わせると、
全 $k$ で $\operatorname{Ku}(V(\sum_iF_i))^{Z_0}\simeq D^b(\prod_iC_i)$ となる
（定理\paperref{main}{mbrd:thm:split-product}）。
二曲線の場合の categorical cover は既に
\cite[Example~4.7]{MBRDBFKFurther}にあり、Hu\cite{TBHu}は幾何的比較を与える。
本稿では全 $k$ への既知定理の反復適用を明記し、
中心作用の coherent な同定と非可換一般点の有限核降下へ接続する。
可換側の積公式そのものを新しい一般定理とは数えない。
明示的な Jacobian tensor 族ではさらに、
$J_i=g_i+\epsilon_i^2\det\operatorname{Hess}(g_i)/24$ を入力の点三次とすると、
復元模型を $\sum_i c_i\lambda_iJ_i$ と取れる
（定理\paperref{main}{mpc:thm:point-model}）。
Kontsevich のグラフ構成\cite[\S\S5.1,6.1--6.3]{NFMKontsevich}の直和性と、
既存の一ブロック全次数恒等式を用いる。
点三次式の既知の計算\cite[Theorem~7.2]{MoriSmith}と、
本稿の同変残余比較への接続を区別する。

この点三次への対応には、中心線の尺度を含む自然な形がある。
$S_i$ の最大可換商の三次部分 $Q_{i,3}$ が直線束で、
中心三次線 $C_i$ と点三次線 $J_i$ からの射がともに可逆である非空開では、
その商を介した同型 $\theta_i:C_i\simeq J_i$ が全次数付き同型と底変更に自然である
（命題\ref{mbpl:prop:line-comparison}、命題\ref{mbpl:prop:nonempty-open}）。
相対 theta 群の枠を選ぶ滑らかな被覆上で
$\theta_i(\Omega_{H,i})=3(a_i-b_i)F_{H,i}$ と書ける。
既知の Hesse 中心式\cite[\S2, p.~3]{GRSklyaninDLL}の非零性と、
分母のない最大可換商の恒等式を用いて、この比較を基本楕円開全体へ一意に延長する
（定理\ref{mbpl:thm:global-line-comparison}）。
相対被覆には Mumford--Shin の theta 群の理論
\cite{MBPLMumfordTheta,MBPLShinWeil}を用い、
普遍被約重なりで降下した後に任意の非被約底へ引き戻す。
従って中心関係直線 $\ell$ を点三次の分離和へ写すことにより、
滑らかな射影三次の stack への射 $\Phi_{\rm pt}$ を得る。
標的の Brauer--Severi 表示を許せば、
$\beta=\Phi_{\rm pt}^*\beta_{\rm cub}$ と $\ell\simeq\ell_{\rm pt}$ を
同時に、しかも自然に実現する（定理\ref{mbpl:thm:commutative-family}）。
この延長は、全 $\OO$-認識自体の開を広げるものではない。
また、二捩れ境界では比較係数が零となるため、同じ線同型を延長できない。

可換分離和そのものについては、滑らかな準射影基底 $S$ 上で
$f_i\in\Sym^3W_i\otimes L$ を共通の係数直線 $L$ に値を取る滑らかな三次とすると、
相対 Orlov 比較\cite[Corollary~3.4]{MBRPBDFIK}と実 dg Koszul 模型から
\[
 \Perf(E_1\times_S\cdots\times_SE_k)
 \simeq\Ku_{\Perf}(V(\textstyle\sum_i f_i)/S)^{Z_0}
\]
を実際の perfect 核で構成する（定理\paperref{main}{mbrp:thm:relative-product}）。
共通の周期直線を含む Hom 複体の外積を直接照合し、
最小 $A_\infty$ 模型に関する同文献の予想は用いない。
核の表現は To\"en\cite[Theorem~8.15]{MBRPToenMorita}を
滑らか proper な相対 scheme に適用して得る。
構成済みの核・逆核・projector を引き戻すことで、
準コンパクト・分離的 noetherian 底への perfect 圏の比較も従う。
底変更の平坦性や底の被約性は要求しない
（系\paperref{main}{mbrp:cor:base-change}）。
この相対可換比較は既知の Orlov--Thom--Sebastiani 法の本設定への相対的な実施であり、
非可換側からの相対核を新たに存在させるものではない。
体上の非可換実表示には、さらに直接の外積法を適用できる。
THCR の多重斉次イデアルを使って全 $k$ の tensor 積の noetherian 性を証明する。
Keeler\cite[Theorem~2.7, Corollary~3.5]{MBNPKeelerMultiple}の二因子定理を
形式的に反復せず、NC-ampleness と多重次数の端に関する帰納を用いる。
各因子の有限 Koszul 解消の曲率付き stabilization を tensor し、
全体の augmentation stabilization と同定する。
Burke--Stevenson\cite[Theorem~6.4]{LP-BS}の window と合わせて
$(\mathcal K_A^{\mathrm{line}})^{Z_0}\simeq\Perf(\prod_iE_i)$ を得る
（定理\paperref{main}{mbnp:thm:product}）。
有限群の双対同変化と既存の coherent 飽和判定により、
$\mathcal K_A^{\mathrm{line}}=\mathcal K_A^{\mathrm{coh}}$ も従う。
ここでは一変数変形や形式比較の存在を仮定しない。
さらに Keeler の相対 Fujita 消滅
\cite[Theorem~1.5, Definition~2.9, erratum Theorem~E2.5]{NFMKeelerAmple}を使うと、
相対次数が全方向で十分大きい線束の消滅を一様に得られる。
この消滅と同じイデアル層の安定化により、全因子数の noetherian 性を
任意の noetherian 基底上で示す。
滑らかな基底では Burke--Stevenson の次数零環に関する仮定も満たされ、
有限 Koszul--Shamash 外積と window を相対化できる。
その結果、順序付き実表示上で
\[
 \mathcal K_{A/S}^{Z_0}\simeq\Perf(E_1\times_S\cdots\times_SE_k/S)
\]
を smooth proper な相対 compact 圏の同値として構成する
（定理\paperref{main}{mbrn:thm:relative-product}）。
非正則な底への引戻しでは、この構成済み perfect 模型を用い、
絶対的な全 coherent 特異圏との一致を要求しない
（系\paperref{main}{mbrn:cor:base-change}）。
可換積側も、相対核の底変更により元の体上で比較できる
（系\paperref{main}{mbrp:cor:comparison-field}）。
通常の Ku 自体は、積圏の双対指標作用による同変化として記述できるが、
積圏そのものとは同一視しない（系\paperref{main}{mbrd:cor:dual-action}）。
非可換側の因子 rotation は
$R_{i,p_i}=T_{\OO_{E_i}}\circ(t_{p_i})_*\circ(-\otimes L_i)$ である。
$p_i$ を表す実線束の tensor により可換側の rotation と共役し、
指定周期の差を三乗根で補正する。
各曲線の次数三の閉点とこの補正から、
順序付き実表示では次数高々 $3^{2k}$ の有限拡大後に
$\Ku_{\mathrm{coh}}(A)\simeq\Ku(X_{\mathrm{pt}})$ を、
中心作用・rotation・指定周期を保って構成できる
（定理\paperref{main}{mbrc:thm:finite-point-comparison}）。
特に代数閉体上の全ての基本領域の点に適用できる。
従って、同じ偏極点曲線の組を持つ幾何点では、
標識を忘れた残余圏の同値類は一定である
（系\paperref{main}{mbrc:cor:point-fiber-constancy}）。
これは比較関手の一意性や相対的な自明化を意味しない。
この体上の比較を、全 $\OO$-基底上の核の大域的 section、
あるいは指定形式核の全 jet の代数化と同一視しない。
この大域問題について、両側の smooth proper な族と指定周期を固定する。
残余圏の周期には引き続き MF/window 比較を要求する。
局所的に rotation を intertwine する
scalar-simple な核の局所型を固定し、負の自己 Ext 消滅と底変更を仮定する。
核の実現 gerbe の類を $\alpha_P$、関係直線の同期を $\theta$ とすると、
rotation 比較の直線束 $L_P$ には
$q_P:L_P^{\otimes3}\simeq\ell_t\otimes\ell_s^\vee$ が付随する。
$q_P(\eta^3)=\theta$ を満たす比較は $\mu_3$-torsor であり、
その類を $\gamma(P,\theta)$ とすると実現条件は
$\alpha_P=0$ と $\gamma(P,\theta)=0$ である
（定理\ref{mbks:thm:kernel-criterion}）。
前者の band は周期を標識しても $\mathbf G_m$ のままである。
この判定は Kummer 理論\cite{NFMStacksKummer}の相対的な適用であり、
積 Picard 補正型では、さらに両障害を因子データで記述できる。
$p_i$ の Picard gerbe の類を $\delta_i$ とすると、外積の band は
$(u_i)\mapsto\prod_i u_i$ なので $\alpha_{\rm Pic}=\sum_i\delta_i$ となる
（命題\ref{mbpd:prop:product-gerbe}）。
因子比較の Hom 直線は
$\lambda_i=(f_i)_*\Omega^1_{E_i/S}$ に標準同型であり、
指定周期から $q_i:\lambda_i^{\otimes3}\simeq\OO_S$ を得る。
$j_i\ne0,1728$ の開では $h_i=c_{6,i}/c_{4,i}$ として
$u_i=q_i^{\otimes2}(h_i^{\otimes3})$ と置けば、
因子周期の類は $\gamma_i=[u_i]$ である
（補題\ref{mbpd:lem:hodge-line}、系\ref{mbpd:cor:unit-phase}）。
この三乗写像そのものも、自然な中心・点三次同期に対して計算できる。
Jacobian の Weierstrass 表示で
$\omega=dx/(2y+a_1x+a_3)$、$v_\omega=2y+a_1x+a_3$ と置けば、
\[
 q_i(\omega_i^{\otimes3})=-v_{\omega_i}(p_i)^{-1},\qquad
 \gamma_i=[v_{\omega_i}(p_i)]
\]
である（定理\ref{mbhp:thm:period-formula}）。
Weierstrass 座標と微分の変換則は古典的である\cite{NFMSilvermanArithmetic}。
本稿では実際の中心三次を全 Hesse パラメータ曲線上で正則化し、
原点での MF 周期と graph--diagonal の余法写像を比較して、
この指定周期の係数と符号を固定する。
二捩れ点での局所付値だけから大域式を推測する議論ではない。
中心作用を保つ Picard 型の比較 stack は
$\mathscr G_{\rm Pic}\times D$、
rotation と指定周期まで保持するものは
$\mathscr G_{\rm Pic}\times\prod_iT_i$ となる
（定理\ref{mbpd:thm:coherent-corrections}）。
従って前者では $\sum_i\delta_i=0$ と全 $\gamma_i$ の一致、
後者では $\sum_i\delta_i=0$ と全 $\gamma_i=0$ が必要十分である。
前者の体上の存在次数は $3^{2k-1}$ に改良される
（系\ref{mbpd:cor:field-degree}）。
この分類では中心作用の標識を保持しており、
その標識を忘れた全比較核の分類とはしていない。
自然な同期でこれらの位相が消えない例も構成する。
Hesse 族 $(a,b,c)=(1+d,1,c)$ では、$d=0$ の点曲線と中心元は非退化だが、
自然な target の方程式は $3dF_H$ となる。
最高 Koszul syzygy の直線を保持した一因子 window をこの境界へ延長し、
正則な auxiliary 方程式 $F_H$ の周期と比較すると、
自然な位相の剰余は $1\in\mathbb Z/3$ である
（系\paperref{main}{mbrn:cor:factor-boundary}、命題\ref{mbpd:prop:boundary-phase}）。
独立な $k$ 変数を使えば、$\alpha_{\rm Pic}=0$ でありながら、
上記 Picard 型の中心作用付き比較に必要な最小拡大次数が $3^{k-1}$、
rotation と指定周期も保持する場合は $3^k$ となる。
これは有限型族の函数体 $\mathbb C(d_1,\ldots,d_k)$ でも成立する
（系\ref{mbpd:cor:phase-independence}）。
Orlov のアーベル多様体に対する
Fourier--Mukai 完全列\cite[Propositions~2.15, 3.3, Theorem~4.14]{MBFCOrlovAbelian}
を使うと、Orlov 像が block 対角である自己同値は因子ごとへ分解できる。
非同種因子では全ての核がこの条件を満たす。
さらに $k\ge3$ では、全双対中心作用との可換性だけで非対角成分を消せる。
実際、異なる二因子の間の成分 $u$ に対し、第三因子を使った作用を取れば
$(r-I)u=0$ となり、$r^2+r+I=0$ から $3u=0$、従って $u=0$ となる
（補題\ref{mbfc:lem:marked-block-rigidity}）。
この整数 Hom 群上の証明は同種性や CM に依存しない。
三次周期が因子間の余分な shift を除き、全 underlying 比較型は
偏極自己同型と rotation・shift の外積から得られる。
体上の scalar 障害はなお $\sum_i\delta_i$ であり、
位相条件は $\gamma_i+\nu_{E_i}(\phi_i)$ の一致、
指定 rotation・周期付きではその全ての消滅となる
（定理\ref{mbfc:thm:full-criterion}）。
ここで $g(f_i)=c_gf_i$ なら $\nu_{E_i}(\phi_i)=[c_g]$ であり、
一般の体上では非零になり得る。
固定 Hesse level の偏極自己同型は既知の十八元であり
\cite[Section~4, pp.~246--247]{NFMArtebaniDolgachev}、
三次式を厳密に保つ lift があるため、この追加位相は消える。
従って上の独立境界族の最小次数 $3^{k-1}$、$3^k$ は、
全ての中心同変核型について鋭い（定理\ref{mbfc:thm:sharp-all-types}）。
三ブロック以上の作用の正規化群が誘導する $H$ の自己同型は、因子置換だけである。
指標の共通反転は、因子周期の式 $[2]=[-2+3n]$ と矛盾して実現されない
（命題\ref{mbfc:prop:normalizer}）。
従って相異なる $j$ の三因子以上、または非同種の二因子では、
中心群の生成元の標識を外して部分群として保つ場合にも判定は変わらない。
二つの因子が同種の場合は異なり、任意の isogeny から Pell 方程式を使って
無限位数の混合 Orlov isometry を作り、実際の coherent 中央化核へ持ち上げる
（命題\ref{mbfc:prop:two-isogenous-mixing}、補題\ref{mbfc:lem:cyclic-centralizer-lift}）。
その各型の定義体、scalar gerbe、周期位相は引き続き算術的な課題である。

中心部分群を保つという条件そのものは、通常の残余圏の同値からは従わない。
二ブロックの滑らかな分離和四様体では、二つの flex 接線から三つの平面
$P_0,P_1,P_2$ を作ると、中心生成元がそれらを巡回する。
平面イデアル層は Kuznetsov 圏の二次元球面対象であり、
その数値類の Mukai 行列は対角成分 $-2$、非対角成分 $2$ となる。
Seidel--Thomas の既知の球面捻り\cite[\S2b, Proposition~2.10]{MBUNSeidelThomas}を
この三つの対象へ適用し、実際に中心を正規化しない自己同値を得る。
二つの捻りの積の冪から中心の無限個の相異なる共役部分群が生じる
（定理\ref{mbun:thm:non-normal}、系\ref{mbun:cor:infinite-orbit}）。
既存の中心同変比較により、この障害は基本領域の二ブロック非可換残余にも移る
（系\ref{mbun:cor:nc-non-normal}）。
従って無標識比較では、中心作用の共役軌道とその算術的降下を別に調べる必要がある。
この非正規性だけで無標識比較の最小次数が小さくなるとは結論しない。
また、一般 scalar 障害を含む上界の鋭さも未解決である。

相対積核の構造移送も次数ごとに比較する。
共通 scalar の自然自明化と中心同変対象の構造射が、
$\Gamma_k$-次数 $\boldsymbol a$ 上でともに
$\prod_i\zeta_i^{a_i}$ となることから、核は射影的表示へ降下し、
さらに中心 rigidification 上で非捩れの $\Perf(\prod_iE_i)$ を与える
（定理\ref{mbic:thm:projective-product}、命題\ref{mbic:prop:rigidified-product}）。
Picard gerbe と位相 torsor の積による比較 stack も、
両側の偏極・周期の twist が相殺するため因子モジュライへ降下する。
既存の全構造対象認識開では全 $\OO$ 射に自然となる
（定理\ref{mbic:thm:kernel-stack}、系\ref{mbic:cor:all-o}）。
核の stack が降下することと、その大域切断が存在することは異なる。
個々の rotation の非捩れ降下も、相対的な rotation 比較からは従わない。
既に全データが置換両立している場合には、連結な順序付き被覆上の
周期 torsor の切断は $\Hom(\mathfrak S_k,\mu_3)=0$ により置換不変となる
（命題\ref{mbks:prop:permutation-phase}）。
この限定的な消滅は核自身の $\mathbf G_m$-gerbe へは適用できない。
Sklyanin の幾何、Heisenberg 群、Zhang twist、同変圏と有限核降下の既知理論と、
本稿の中心商・第三 Veronese・同変一般点復元への適用を区別する。
これらの組合せの新規性について網羅的な確認は未完了である。
量子化論文では周囲代数を$A$、中心商を$B$と記すのに対し、
本稿では周囲を$B$、中心商を$A$と記す。
中心元の規格化と中心イデアルを保つかどうかは、引用ごとに区別する。

\subsection{可換復元の入力と先行研究}

可換復元稿の一般中心三次変形の基礎定理と、その一ブロック型への適用は、
論理的にも対象範囲としても異なる。完備noetherian局所底での
有限rank graded MF の同値は定理\paperref{main}{cft:thm:complete}、
$\mathbb C[h]_{(h)}$ 上の次数別有限自由な中心三次変形に対する
全coherent残余圏と可換Kuznetsov成分の一般点同値は
定理\paperref{main}{gdr:thm:generic}による。後者の可換方程式は元の関数体にあるが、
比較核を実現するためには有限拡大を許す。その後で初めて、
良い一ブロック型の$\mathcal O$-marked圏の相対構造と
内在的な残余三次$X_{\mathrm{res}}$の復元・降下を扱う。
一般中心三次変形の定理を、一ブロック表示への存在定理に読み替えない。
以下で用いる可換復元の入力は、可換復元稿の完備MF移送・相対window・
一般点降下・coherent飽和の明示的な定理とその仮定である。
適用する対象の表示、底の正則性、指定周期、飽和性の条件を各主張で保持する。

関連する公刊研究との関係は次の通りである。
Etingof--Ginzburg \cite{NFMEG} は三変数 potential の PBW 変形と中心元を構成し、
Mori--Smith \cite{MoriSmith} は三生成元三二次関係の Calabi--Yau 代数を分類する。
Behrend--Noohi \cite[Sections~2.3, 3.1.2]{NFMBehrendNoohi} は
次数スカラーによる rigidification、重み付き代数の gerbe、
determinant character による中立化、および偏極付き Grothendieck 圏を扱う。
これらの一般的方法を新規とは数えず、本稿では全 $\mathcal O$-同値の復元を用いて
shift の標識を外した局所一ブロック型の圏 stack と同定し、
階数 $(n-1,3)$ に依存する、全ての非空開上の Brauer 障害を決定する。
Berger--Taillefer \cite{NFMBeTa} の potential 解消と PBW--CY 定理も入力とし、
本稿では任意個の中心差の Koszul 解消と有限射影模型によって、
非平坦・非被約な底変更にも両立する相対双対と trace を構成する。
体上の CY 性や PBW 定理自体を新規成果としては数えず、
$n$ に依存する全同型問題や torsor の降下と区別する。
Le Bruyn--Smith--Van den Bergh \cite{LBSV} の一点中心拡大における
楕円曲線と八点、および二次関係の復元に対して、本稿は任意個数の中心変数を扱い、
有限平坦な八重被覆と高次元の二十四次完全交叉を得る。
さらに Rains の中心公式 \cite[\S9, p.~1414]{NFMEG} の境界評価から
中心三次の復元を強め、相対 trace と尺度付き Hessian による
Poisson 対への射を構成する。この射と既存量子化の単純な逆写像、
および残余圏からの可換復元とは区別する。
形式量子化の中心と Poisson 中心の同型は
Etingof--Ginzburg\cite[Proposition~8.1.1, equation~(8.1.3)]{NFMEG}
にも既に現れ、その証明では unimodularity から相対形式 Poisson tensor を
Jacobian potential で表している（Section~8.1, pp.~1407--1408）。
本稿で扱うのは、有限の非零パラメータでの三次部分に対する
trace と Hessian の明示的処方、その任意底変更、全同値に関する降下と逆構成である。
Dito\cite[Section~2, equation~(3); Section~3, Theorem~2]{NFMDito} の一般 Poisson
星積の二次項と Weyl gauge、および Shoikhet\cite[Sections~1.3, 2.3.1]{NFMShoikhet}
の内向き wheel の消滅を用いて、固定入力の残余 potential の二次計算を
全混合係数へ広げる。星積や wheel 消滅自体は既知結果であり、
本稿では trace--Hessian 処方との差が全次元で非零となる開条件を調べる。
二次代数と二次 Poisson 構造の形式対応は既に
Kontsevich\cite[Section~2.3, Proposition~6]{NFMKontsevichAlg}にある。
また点曲線の Hessian 公式は Mori--Smith\cite[Theorems~1.2, 7.2]{MoriSmith}の結果である。
本稿はこれらを既知の入力として、混合係数の trace を加えた代数的補正を
O-marked moduli 上へ降下し、全混合係数での一般次数九と残余比較への適合を調べる。
さらに Lindberg--Pym\cite[Section~1.6; Theorem~6.9 and Remark~6.10]{NFMLindbergPym}
の量子トーラスの交換比 $q=e^\hbar$ を既知の入力として、
生成元固定の逆 tensor の五次項を三変数三次の共変式から決定する。
全混合項の定数評価には
Banks--Panzer--Pym\cite[Corollary~1.3; Section~6]{NFMBanksPanzerPym}の
harmonic graph weights と、Buring\cite[Sections~3.4--3.6]{NFMBuringGraphs}の
公開データを用いる。線型入力と二次 tensor に有効な73項を抽出し、
元の weight と多重度、生成元固定再帰を有理数として検証する。
既知の graph 積分値自体を新規成果とは数えず、
本稿の貢献は混合項の一様な公式、残余三次の四次比較類とその非消滅、
および全同値へ降下する追加補正への適用である。
形式対応の存在、点曲線公式、量子トーラスの交換比そのものを新規成果とは数えない。
次数計算に用いる Hesse 還元は、純 Hessian 写像に対する
Ciliberto--Ottaviani\cite[Theorem~5.7]{NFMCilibertoOttaviani}の議論にもあり、
四次の $S$ は古典的な Aronhold 不変式\cite{NFMOttavianiAronhold}である。
$I,J$ も Hessian 写像と Aronhold 四次式の方向微分で表示できる
（補題\ref{lem:nfm-classical-mixed-covariants}）。
本稿ではこれらを一般混合係数と商 stack の次数、および固定逆 tensor の係数へ適用する。
Mayanskiy \cite{Mayanskiy} の三様体における Schouten 方程式の計算に対して、
本稿は高次元での Jacobi イデアル自体を比較する。
affine 超曲面の Poisson 延長を gradient の Koszul 複体で記述する手法には
Fraenkel \cite[Proposition~4.1, Theorem~4.3]{Fraenkel} の先行研究がある。
また、一般の Lie 代数 law の Jacobi スキームに冪零元が現れることは
Carles \cite[Section~8, Lemma~8.35]{Carles} および
Carles--M\'arquez \cite[Section~3]{CarlesMarquez} に既知である。
本稿の Jacobi 係数写像は Landsberg--Ottaviani の Young flattening
\cite{LandsbergOttavianiYoung}、Gesmundo--Landsberg の Koszul flattening
\cite[Proposition~6.1]{GesmundoLandsbergPartials}に同定される。
従って普遍階数上界自体は既知であり、本稿は座標による再証明を与える。
apolar Gorenstein 代数の Betti 数を計算した Boij \cite[Remark~6]{NFMBoijBetti}
との関係も明示し、$N=9,10$ の二次階数の達成を新規とは数えない。
Jacobi 節で扱う中心的対象は、三次超曲面に付随する Hessian--Jacobi イデアルと
Grassmann 錐のスキーム論的な差、次数三の乗法が決める長さ330の平方零構造、
および特殊三次の非零点の局所構造である。
Hessian と apolar 積、Thom--Sebastiani 分解の関係については
Bricalli--Favale--Pirola \cite[Proposition~2.1, Theorem~3.4]{NFMBFFHessian}を参照し、
本稿の横断切片と Jacobi スキームの retract の主張と区別する。
Lin--Zhang \cite{NFLZSerre}、Pirozhkov \cite{NFPirozhkov}、
Huybrechts--Rennemo \cite{NFHuybrechtsRennemo} の既知の圏論的不変量と Torelli 理論は、
本稿の残余変形の入力である。

残余三次の勾配空間による一意性は、Wang\cite[Theorem~3.3]{NFMWangMilnor}と
Hwang\cite[Theorems~1.5, 1.6]{NFMHwangSymmetrizer}の既知結果に基づく。
本稿では階数一の積分核を相対的に作り、O-同値の線束 twist による重みを追跡して、
Brauer--Severi ambient と Poisson 直線への射を降下する。
perfect な対角・核を有限係数へ広げる一般原理と同値 locus の開性には
To\"en--Vaqui\'e\cite[Lemmas~2.10, 3.25]{NFMToenVaquieModuli}を用いる。
これらの一般原理、体上の三次の存在、既知の量子トーラス交換比を新規とは数えず、
相対的な残余射・中心作用の保存・有界次数の関係 locus による
very general な固定 tensor の超越性への適用を区別する。

点曲線と残余 active 曲線の比較には、既存の split Hesse 同定を
active 再帰の制限と係数空間の支配性によって一般混合係数へ広げる。
相対随伴による尺度線束と Hodge 線束の双対性、および
Fulton--Olsson\cite[Section~1; proof of Corollary~2.4]{NFMFultonOlsson}に記載された
判別式と楕円 moduli の $\mu_2$-gerbe は既知の幾何である。
本稿では点曲線の量子化から定まる微分をこの幾何と組み合わせ、
同じ残余三次を標的とする符号商 Poisson 射の支配性・一般有限性を証明する。
平面三次の十八個の自己同型や二つの Jacobian 同型そのものを新規とは数えない。
また Fisher\cite[Theorem~1.1, Lemmas~2.2, 2.6]{NFMFisherTernary}の
三変数三次の不変式環と微分付き軌道分類を用いる。
本稿では active 不変量への幾何学的一般ファイバーが整であることを混合係数も含めて示し、
九次の中間被覆を取り出す。さらに $h$ の全次数にわたる混合重みの有限性と
trace の多項式性から混合係数の多項式逆を作り、共通の三捩れ安定化群の比較により
残余射の一般次数を十八と決定する。
幾何学的一般ファイバーの整性に基づく体論
\cite{NFMStacksRegularFields}を用い、符号二重被覆の類を active 不変量の体へ降ろす。
古典的不変式論や正則な体拡大の一般論自体とは区別する。
さらに Hesse 係数体への二次類の単射性と、三捩れ対称な混合族を用いて、
符号類の消滅を証明する。Benoist\cite[Theorems~1.6--1.7]{NFMBenoistSeparation}による
滑らかな偏極三次の分離性・DM 性から、Poisson 符号商標的の有限対角を得る。
有限表示の係数降下\cite{NFMStacksStackLimits}と正規環への整性によって
指定した一般点での同型を延ばし、混合三角多項式から得る近傍延長を降下する。
正規な底上の有限 étale 被覆の一般ファイバーへの完全忠実性
\cite{NFMStacksEtaleNormal}と有理三捩れ点による twist の排除を組み合わせる。
これらの一般的な降下原理自体は既知であり、残余射への適用と区別する。
比較核の scalar 障害には、Knudsen--Mumford\cite{NFMKnudsenMumfordDet}の determinant と
Lieblich\cite[Sections~2.2.4--2.2.5]{NFMLieblichTwisted}の twisted sheaf の重みの議論を使い、
可換復元の入力における stabilization の Euler 数から全次元の明示的な index bound を取り出す。
さらに次数一の中心正則列に対する graded purity と active 曲線への制限により、
最高次 cohomology が線束となる型で次元に依存しない階数三の証明書を作る。
三変数の $a=0$ 比較は Orlov\cite[Theorem~2.5(iii)]{NFMOrlovSingularities}と、
正則係数環への拡張である Burke--Stevenson\cite[Theorem~6.4(3)]{NFMBurkeStevenson}を入力とする。
window と内部 shift の規約は Galinat\cite[Section~1.1, Lemma~1.2; proof of Lemma~2.7]{NFMGalinat}
に合わせる。これらの同値や楕円曲線上の spherical twist 自体を新規とは数えない。
determinant による Brauer 障害の消去原理自体、および極大部分体による分裂
\cite{NFMStacksSplitting}は既知結果である。
Picard stack の gerbe 構造\cite{NFMStacksPicardGerbe}と
重み零降下\cite[Proposition~4.10]{NFMBerghSchnurer}も既知である。
本稿ではこれらを最高次線束へ適用し、三倍の Picard 類を表す線束を
determinant で明示的に作り、選択済み核を相対的な三分点被覆上で規格化する。
楕円曲線の三倍写像の次数九\cite{NFMStacksMultiplication}自体は新しい結果ではない。
右 twisted homogeneous coordinate ring の shift は
Stafford--Van den Bergh\cite[\S3, equation~(3.1)]{NFMStaffordVdB}の規約を用いる。
この shift と spherical twist の共役を計算し、
指定した中心三次を持つ active 商に対し、相対 THCR と window1 の比較を構成する。
相対 ample pair は Keeler\cite[Proposition~7.2, Lemma~7.6]{NFMKeelerAmple}の枠組みを用い、
底変更は各次数の perfect cohomology\cite{NFMStacksPerfectBasechange}と
MF の余核図式から検証する。これによって固定入力の下で、
最高次線束の Picard 類を点の平行移動へ同定する型の非空性を得る。
体上の同値の存在だけから相対同値や底変更を推論してはいない。

genus-one 曲線に付随する Jacobian 上の Azumaya 代数と、その点における
Brauer 評価は既知である
\cite[Theorems~1.1--1.2, Corollary~1.4]{NFMFisherGenusOneAlgebras}。
Ciperiani--Krashen\cite[Theorems~3.1, 3.7]{NFMCiperianiKrashen}は
Picard 表現障害を Tate pairing と同定し、有限 torsion の Weil cup 積で表示する。
本稿ではこの既知公式を flex torsor、平面同型 torsor、三分点 torsor へ適用し、
二つの active 曲線間の障害の差を具体化する。
さらに一般の $j$ を持つ Hesse 型曲線上で接線の norm を直接計算し、
非零の index 三 Picard 障害を求める。巡回代数や norm 判定\cite{NFMMilneCFT}自体は既知である。
Matsuno\cite[\S2.4, Lemma~2.16, Theorem~3.3, \S4]{NFMMatsunoElliptic}の
幾何的 Sklyanin 表示と cyclic potential、および
Mori--Smith\cite[Proposition~2.11, Theorem~4.2(1)]{MoriSmith}の
基底変更と正則性判定を用い、この例を元の体上で勾配代数に実現する。
ここで本稿が確認するのは、原点や Picard 類の線束代表を仮定しない potential の降下と、
非零 Brauer 類を保持する支配的な全係数表示族である。
これにより一般 split 表示の点障害の正確な index 三を得る。
さらに Fisher\cite[proof of Lemma~2.3, \S9.5]{NFMFisherTernary}の
平面三次の三捩れ作用と Heisenberg 群を使い、
その全 ambient 上の中央化群を計算する。
非零の自明中心部分を伴う表現では中央化群は split な $\operatorname{GL}_{n-1}$ であり、
active 平面上では恒等に作用する。
有限対角延長と有限 torsor の混合定数降下を組み合わせると、
一般係数体で半古典平面同型 torsor が分裂し、残余側の Picard 障害も正確に index 三となる。
Heisenberg 群自体や Hilbert 90 を新規結果とは数えず、
本稿で証明するのはこの full ambient と残余射への適用である。
形式平面比較だけでは同じ結論が出ないことは、明示的な巡回 twist の相殺例で区別する。

楕円曲線上の Fourier--Mukai 同値の数値作用とその核は
Hille--Van den Bergh\cite[Proposition~4.2, Corollary~4.3, Theorem~5.3]{NFMHilleVdBFM}
に記述されている。本稿ではこの graph 認識の議論を実際の核に適用し、
三次 rotation の数値 centralizer の六元計算と合わせて、
原点を仮定しない体上の全 intertwiner の存在判定を与える。
相対的には、非被約底を許す normalized kernel stack と Picard gerbe を同定する。
underlying kernel の分類と、rotation 同型自体および指定周期を標識した際の
追加の $\mu_3$-torsor を区別する。

三・四次元の full 核の同型類については、
Liu\cite[Theorem~4.13]{NFMLiuCubicThreefoldAut}と
Huybrechts\cite[Theorem~1.2(i)]{NFMHuybrechtsK3Cubic}の既知の自己同値分類を用いる。
前者は滑らかな複素三次三次元多様体、後者は very general な複素三次四次元多様体を扱う。
一般三次の幾何自己同型の自明性\cite[Theorems~1.1, 1.5]{NFMPoonenAutomorphisms}と合わせ、
両次元の幾何学的一般点では $\rho^{-1}[1]$ が全群を生成する。
これらは関手の同型類の分類なので、核への適用には追加の一意性が要る。
本稿では projected point の simple 性と
Canonaco--Stellari\cite[Corollary~4.4]{NFMCanonacoStellariKernels}の cohomology 一意性から
三次多様体の projector 核の一意性を証明する。
その後、定数無限巡回群の連続 Galois cocycle が零であることを用いて全 full 類を降下し、
既存の Picard 障害と接続して最小比較次数三を得る。
自己同値群の既知分類自体を新規成果とはせず、very general な分類を Zariski 開の主張に置き換えない。
指定周期を加えた比較では、同じ整数冪による移送が周期位相も保つことを示す。
中心商の conormal 同定を用いて方程式生成元の相対尺度の変換則を追い、
位相を消す規格化と、独立な標的尺度を付した一般族での正確な最小次数九を得る。
さらに full tensor と半古典平面同型から方程式生成元を三乗倍まで規格化し、
周期類 $\gamma_{\mathrm{sc}}$ を定める。
この類は形式比較体へ移すと零になるが、後述の Hesse 境界計算により
元の係数体上では非自明であることが分かる。
Murfet\cite[Theorem~6.2, equation~(6.6)]{NFMMurfetResidue}の residue trace は
可換 MF 側での計算法を与える一方、source trace と CF 移送の scalar 比較を別に要する。

高次元では、rotation 中央化群 $G_\rho$ の定数中心部分
$\langle T\rangle$ を除いても各 class の Galois 安定化群が変わらないことを示す。
奇数次元の数値格子の計算から、残る群は
$G_{\rho,0}/\langle[2]\rangle$ と同定される。
これは点型 class の選択を具体化する簡約であり、その存在の証明とは区別する。
本稿ではさらに、最高次線束 $M_P$ の右随伴像から原点像を復元する条件を課す。
形式恒等分枝における Koszul 計算と proper な導来 Nakayama によって、この強い型の非空性を示す。
三連続の原点対象の生成性と $HH^1=0$、
To\"en--Vaqui\'e\cite[Proposition~3.20, Corollary~3.24]{NFMToenVaquieModuli}の
bounded な対象 moduli と自己同値 locus の開性から、その class を動かす群 $H_p$ は有限となる。
従って無限な自己同値群の全分類を先に解かずに、有限 torsor と Picard gerbe を分離できる。
さらに本稿では、全 $n\ge3$ の一般一ブロック族で $H_p=1$ を証明する。
中央三次商の双加群自己交叉から通常 Ext と Tate Ext を比較し、
双次数による切断で通常 Ext を回復する。
Lu--Palmieri--Wu--Zhang\cite[Theorem~2.4, Corollary~2.5]{NFMLPWZKoszul}の
Adams Koszul 双対と、Braun--Chuang--Lazarev
\cite[Proposition~3.17, Theorems~3.18, 5.3]{NFMBCLLocalization}の
derived localization の忠実性を組み合わせ、原点・rotation を保つ核を
$\operatorname{Aut}_{\mathrm{gr}}(B,k\Omega)/\mathbf G_m$ へ単射で移す。
Keller\cite[Sections~9.3--9.4]{NFMKellerOrbit}の enhanced orbit の関手性を用い、
内部次数を保つ covering category へ戻すことで、対象像だけでなく核の class を検出する。
一般表示の射影安定化子は自明なので、強化型の class は全次元で一意となる。
Koszul 双対や局所化そのものの新規性は主張せず、ここでの cutoff と
比較核の算術的 class 選択への適用を区別する。

同じ復元から、指定周期の位相を有限な低次積へ帰着する。
Picard gerbe 上の $Z_p=q_Y^R\mathcal M[1]$ の Adams graded orbit は
重み零なので元の体へ降り、その二次積が周囲代数 $B_p$、
指定周期と三次積が $\Omega_p\in(B_p)_3$ を与える。
Lu--Palmieri--Wu--Zhang\cite[Corollary~B]{NFMLPWZExt}の関係復元を用い、
三次 tensor の gauge 差を二次関係 ideal で除いた上で、
$f:(B_p,K\Omega_p)\simeq(B,K\Omega)$ に対する
$\gamma_{\mathrm{sc},n}=[f(\Omega_p)/\Omega]$ を証明する。
元の体上の実際の比較核を仮定せず、標的の方程式の尺度と
source の自己交叉周期の符号を同じ式で保持する点が、この適用に必要である。
良い integral 表示と原始的な三次関係が閉 fibre にもある場合には
非分岐性も従う。一方、以下の境界では二次関係を飽和した格子内で
三次関係が境界パラメータの一倍を含み、この十分条件を満たさない。

具体的な $q=2$ の Hesse 境界では点曲線は split $I_3$ に退化する。
Silverman\cite[Chapter~XI, Proposition~9.2]{NFMSilvermanArithmetic}の
既知の Kummer 評価公式へ規格化した接線比を代入し、
三分点被覆が三つの同じ完全分岐三次体に分かれることを計算する。
一方、点線束 $\mathcal O(P-O)$ の三行三列線型表示を、完備化前の trait 上の
有理式として構成する。点尺度を保つ中心座標の変更で、実際の一般残余三次の
平坦な模型を取り、その上の右射影対象を有限 mutation で構成する。
選択した模型の特殊 fibre は特異でもよい。
Eisenbud\cite[Theorem~7.2]{NFMEisenbudCI}の既知の構成による有限 MF を明示し、
Burke--Stevenson\cite[Theorem~6.4(1)]{NFMBurkeStevenson}の相対窓を指定して
右射影対象と同定する。この窓は特異な射影 fibre にも適用され、
内部次数と同じ方程式の MF 周期を保つ。
特殊 fibre の active MF と中心 Koszul 因子の分離により、
全 $n\ge3$、全ての整数 $i,j$ について
\[
 \operatorname{rank}_R H^i\operatorname{RHom}(Z_R,\rho^jZ_R)
 =\binom{n+2}{3i+2j}
\]
となり、これらの Hom は自由で閉・一般 fibre への底変更と可換である。
範囲外の二項係数は零とする。従って指定 target 周期は primitive であり、
必要な $m_2,m_3$ を integral に取れる。

しかし特殊 fibre の二次積の階数は $\binom{n+2}{2}-3$ に落ち、
三次 class も零となる。実際の行列の一次変形を計算すると、
失われる三つの二次関係の Smith 指数は全て一である。
二次関係を飽和して得る代数は全次数で平坦となり、
その active 特殊 fibre は三角形型 Sklyanin 代数 $S(0,-2,-p)$ となる。
指定周期から回復される三次 class の初項は、境界パラメータ $z=q-2$ と
非零中央三次の積である。これだけで一般 fibre の表示同型の integral 性は従わないが、
cyclic potential の alternating 成分から滑らかな potential 三次を回復できる。
Benoist\cite[Theorems~1.6(ii), 1.7, Lemma~2.2(iii)]{NFMBenoistSeparation}の
既知の分離性をこの三次へ適用すると、指定した一般 active 射影同型が延びる。
残る自己同型の中央三次への作用が三乗倍であることを用いて、
全 $n\ge3$ の完備局所位相を $[D_R]$ と決定する
（命題\ref{prop:nfm-boundary-period-local-kummer}）。
この結果は滑らかな特殊残余三次や full 比較核の境界延長を必要としない。
楕円 Kummer 理論、$I_3$ 成分群、相対 Orlov 窓、Shamash 構成や
滑らかな三次の分離性そのものを新規成果とせず、
ここでの integral 行列・低次積・規格化周期への適用を区別する。

五次元の安定性条件には Peize Liu\cite[Theorem~1.1, Remark~5.6]{NFMPeizeLiuFivefold}の
Serre 不変な条件の構成があるが、同論文はその軌道の一意性を仮定していない。
本稿では同論文の Clifford rotation を普遍平面上の相対核として構成し、
$HH^1=0$ と Fano 曲面の連結性から、その核 class の平面非依存性を示す。
Clifford 比較の入力は Kuznetsov\cite[Theorem~4.2]{NFMKuznetsovQuadric}と
Liu\cite[Theorem~2.8, Corollary~3.10]{NFMPeizeLiuFivefold}、
曲面の連結性は Collino\cite[Propositions~1.8, 1.10]{NFMCollinoFivefold}による。
二捩れの実現 gerbe を、既知の奇数 Pl\"ucker degree
$2835$\cite[Example~4.4]{NFMCilibertoZaidenbergFano}によって中立化し、
元の体上の半 rotation の実核を得る。
さらに spinor 列を相対随伴の unit と同定し、相対・絶対 coevaluation を比較すると、
$O\simeq T=\rho^{-1}[1]$ という perfect 核の自然同型を得る。
従って $\epsilon=OT^{-1}\simeq\operatorname{id}$ である。
この同定は平面が定義された任意の滑らかな五次元三次に成立し、
数値作用や安定性作用の忠実性を前提としない。

非分裂表示では、実際の次数三線束と偏極類の区別が指数に現れる。
theta group による divisor-class pair の降下
\cite[Proposition~2.2, Lemma~2.4]{NFMCremonaFisherDescentAlgebra}と
Tate--Weil 公式\cite[Theorem~3.7]{NFMCiperianiKrashen}を用いて、
次数零 Picard 障害が period 三・index 九となる genus-one 曲線を構成する。
Brauer--Severi 平面と三捩れの Heisenberg 表現の関係は
Saltman\cite[Sections~3--4, Theorem~5]{NFMSaltmanSeveri}でも扱われる。
本稿ではこの既知の構成原理を、混合係数と中心三次 torsor の twist によって
全 $n\equiv1\pmod3$ の良い一ブロック成分の任意の非空開へ実現し、
点型 full 比較の九整除性と普遍 Picard gerbe の最小正階数九へ接続する。
表示安定化子が自明な共通稠密開内では、全該当次元で
随伴復元点型の唯一の class により最小実現次数九を達成する。
さらに非分裂 rotation 自身の実現 gerbe が偏極障害 $\beta$ であることを示し、
$\rho^rP$ の障害を $\alpha+r\beta$ と計算する。
次数三 Picard 類の標準的な押し下げと低次元自己同値分類を組み合わせると、
全 full 比較の最小次数は $\alpha\in\langle\beta\rangle$ に応じて一または三となる。
従って四次元の指数九族でも全型の最小次数は三であり、点型を保つ比較との相違が現れる。
曲線の period--index の相違そのもの、既知の degree $2835$、
Clifford 半 rotation の存在自体を本稿の新規成果として数えていない。

さらに Wu\cite[Theorems~1.1, 1.3, Proposition~7.2]{NFMWuStability}と
Aravena--Liu--Tevelev\cite[Theorem~1.2]{NFMAravenaLiuTevelevGepner}は
超曲面圏の Gepner 安定性条件に関する一般的な構成を述べている。
これらの存在定理と、自己同値が特定の成分を保つこと、
その作用の忠実性および係数体上の Galois 降下は異なる問題である。
本稿の有限群への還元は、これらの安定性条件の存在・一意性を証明入力に使わない。
また Iliev--Manivel\cite[Proposition~3.1, Corollary~5.3]{NFMIlievManivelSevenfold}の
既知の球面対象を用いると、一般三次七様体の full 自己同値群は
$\mathbb Z^2$ を含み、rotation 中央化群の指数は無限になる。
球面対象の構成自体は既知入力である。この帰結は、低次元の全群の巡回性を
全高次元へ広げる方法に対する障害を明示する。
Torelli 型の対象復元や Fourier--Mukai partner 数の有限性
\cite{NFPirozhkov,NFMLinRennemoZhangWeighted,NFMLinZengZhangSingularity,NFMLinZhangPartners}を、
核の同型類の一意性または降下と同一視しない。

全次元の三分点被覆については、普遍 Hesse 曲線の真の一般点から translation 群を得る。
Fisher の三捩れ安定化群を $g$ と $J_g$ で同一視して線型 monodromy を調べ、
標準 split 係数体上の affine monodromy が全 $\operatorname{ASL}_2(\mathbf F_3)$ であることを示す。
三倍写像の次数や古典的 Hesse 群自体と、この係数体上の連結性・分裂次数の決定を区別する。

既知の公式からの帰結と、本稿に証明を与える拡張は以下で区別する。
これらの拡張について、独立な新規性の網羅的確認はなお完了していない。

\subsection{中心復元と構造対象付きモジュライ}
以下 $n\ge3$、$N=n+2$、$r=n-1$ とする。
生成元空間を $V$、中心一次空間を $U\subset V$、$\dim U=r$ と書く。
点の空間は $W=V^*$ であり、$\PP(W)$ 上の三次式は
$\Sym^3W^*=\Sym^3V$ の元である。
周囲の二次代数を $B$、中心三次商を $A=B/(w)$ と書く。

\begin{maintheorem}[中心復元と局所一ブロック型モジュライ]
可換 $\C$-代数上の次数一生成代数 $A,A'$ が、それぞれ
中心多項式環 $C=\Sym U,C'=\Sym U'$ 上平坦、
$\rank U=\rank U'\ge2$ であり、
普遍的に $Z(A)_d=\Sym^dU$ と $Z(A')_d=\Sym^dU'$ を満たすとする。
このとき $d$-Veronese 間の次数付き同型は fpqc 局所で $A\simeq A'$ へ持ち上がり、
持上げ層は $\mu_d$-torsor である。

一ブロック PBW・楕円中心条件を満たす非空開上で、
次数スカラーを除いた次数付き周囲代数の同型問題は、
有限安定化子開に制限すると滑らかな既約有限型 DM stack $\mathfrak B_n$ となる。
中心三次直線の stack $\mathfrak C_n$ からの射
$\tau_n:\mathfrak C_n\to\mathfrak B_n$ は、
階数 $s_n$ のベクトル束の加法群の torsor であり、
\[
 \dim\mathfrak B_n=b_n=\frac{n^2+7n-4}{2},\qquad
 s_n=\binom{n+1}{3},\qquad
 \dim\mathfrak C_n=b_n+s_n=\binom{n+2}{3}+3n-2.
\]
普遍相対双対と全同値比較の良い開に制限すると、
$\mathfrak C_n^{\mathrm{good}}$ は fpqc 局所的にこれらの表示を持つ
$\mathcal O$-marked Grothendieck Abel 圏族の stack
$\mathfrak M_n^{\mathrm{loc}}$ に同値である。
大域の次数付き表示を妨げる gerbe の類 $\beta$ は
$\gcd(n-1,3)\beta=0$ を満たす。
$n\equiv1\pmod3$ では、その制限は任意の非空開上で位数がちょうど三である。
\end{maintheorem}

正確な普遍中心条件と証明は定理\ref{thm:nfm-central-veronese}に、
表示空間と全次数付き同型の比較は命題\ref{prop:nfm-graded-presentation}に、
有限安定化子と torsor の証明は定理\ref{thm:nfm-ambient-dimension}、
\ref{thm:nfm-central-torsor}に置く。
相対 Serre 関手が与える $d=r$ にこの復元を適用するため、
楕円自己同型 $\sigma$ には
$\sigma^j\ne1$ $(1\le j\le\max(3,r))$ を課す。
中心スカラーの重みを考慮した束
$\mathcal Hom(\mathcal L,\Sym^3\mathcal U)$ が商 stack 上に降下する。
$\mathcal U$ や中心三次束そのものの降下を主張するものではない。

\begin{center}
\begin{tabular}{r|r|r|r|r}
$n$&可換三次の次元 $M_n$&周囲 $b_n$&中心三次 $s_n$&対 $c_n$\\\hline
3&10&13&4&17\\
4&20&20&10&30\\
5&35&28&20&48\\
6&56&37&35&72\\
7&84&47&56&103\\
8&120&58&84&142
\end{tabular}
\end{center}
ここで $M_n=\binom{n+4}{3}-(n+2)^2=\binom{n+2}{3}$ である。
$c_n-M_n=3n-2$ という数値一致から、可換モジュライへの射の存在や
そのファイバーの次元を結論してはいない。
良い相対 PBW 表示チャートでは、定理\ref{thm:nfm-relative-ambient-cy}により
真正な双加群双対 $\Theta_{B/R}\simeq B(N)[-N]$ を構成する。
定理\ref{thm:nfm-relative-line-cohomology}と\ref{thm:nfm-relative-serre}は、
任意底変更後の compact generation、相対 properness、相対 Serre 関手
$S=(1-n)[n]$ とその trace を与える。
固定表示模型間の係数作用と全 Hom の fpqc 降下を構成することにより、
定理\ref{thm:nfm-fixed-model-comparison}は
$\underline{\operatorname{Isom}}_{\mathrm{gr}}/\mathbf G_m
\simeq\operatorname{Eq}_{\mathcal O}$ を与える。
大域的な次数付き代表の障害は $\operatorname{Pic}(R)[n-1]$ の線束であり、
次数一加群がともに自由なら $\operatorname{Pic}(R)[\gcd(n-1,3)]$ に属する。
$n\equiv1\pmod3$ では、滑らかな affine 楕円曲線上で非零障害の例も構成する
（命題\ref{prop:nfm-nontrivial-three-torsion}）。
切断飽和性自体は中心正則対によって全整数次数で証明する
（補題\ref{lem:nfm-central-saturation}）。
定理\ref{thm:nfm-local-moduli}では、係数変更を圏内の係数作用として定義し、
表示 gerbe 上の重み付き次数加群の tails により対象の効果化を構成する。
表示 gerbe の $H^2$ の類と、指定同値の上の $H^1$ の線束障害は異なる。
定理\ref{thm:nfm-brauer-dense-open}は、$n\equiv1\pmod3$ では
どの非空開にも $\C(s,t)$ 上で非零 Brauer 障害を持つ対象が存在することを示す。
全ての内在的成分の一般点に局所一ブロック表示がある、という分類は主張しない。

\subsection{相対点対応と Poisson 対への射}
\begin{maintheorem}[相対点スキーム、中心復元、Poisson 対]
一ブロック模型の楕円境界、固定点なしの shift、および中心三次の正則切断条件の下で、
\[
 \operatorname{Pt}(A)=E\amalg D,\qquad
 D\subset\PP^{n+1}\text{ は型 }(2,2,2,3)\text{ の完全交叉},
\]
\[
 \dim D=n-3,\qquad\deg D=24,\qquad\omega_D=\OO_D(7-n).
\]
三つの偏微分で定まる $C$ は $\PP^{n-2}$ への次数八の有限平坦被覆である。
$n=4$ で $D$ が滑らかなら、$D$ は種数 $37$ の連結曲線である。
これらは任意底変更と両立する相対完全交叉と八重被覆を与え、
点商の関手、shift のグラフ、中心射影は
$\mathfrak M_n^{\mathrm{loc}}$ の良い開上へ降下する。
ambient は一般には Brauer--Severi 束であり、上の双対化層の式は
各ファイバーでの式である。

非空開上では、埋込み点対応から $B$ と中心直線 $[w]$、従って $A$ を復元できる。
さらに尺度込みの周囲 Poisson tensor $\Pi$ に対し、
中心三次束と古典三次 Casimir 束の間に、純中心三次を固定する同変同型
\[
 \mathfrak d_\Pi:Z(B)_3\xrightarrow{\sim}
 \mathcal R_\Pi=\C F\oplus\Sym^3U
\]
を構成できる。この同型は任意の基底変更と両立する。
従って滑らかな三次を与える非空開
$\mathfrak M_n^{\mathrm{Pois}}\subset\mathfrak M_n^{\mathrm{loc}}$ 上で
\[
 (\mathcal C,\mathcal O)\longmapsto
 \bigl(V(\mathfrak d_\Pi(w)),\Pi|_{V(\mathfrak d_\Pi(w))}\bigr)
\]
が、偏極類を持つ三次 Poisson 対への射を定める。
\end{maintheorem}

相対点関手と束の正確な降下は定理\ref{thm:relative-point-nfold}、
\ref{thm:point-omarked-descent}に、中心直線の復元は
定理\ref{thm:point-central-recovery}に置く。
Poisson 射の標的は、斉次 quadratic Poisson tensor と滑らかな三次 Casimir の
局所閉スキームの $\PGL(V)$ 商として具体的に構成する
（定理\ref{thm:point-trace-hessian}）。
定理\ref{thm:poisson-strong-lift}により affine lift は実は Poisson 対から一意に回復する。
従って追加の lift 標識なしに、Brauer--Severi 周囲空間と偏極類を持つ対を扱える。
$n\equiv1\pmod3$ で実際の相対線束 $\OO(1)$ を標識に加えるには、
上記の Brauer 障害への対処が必要である。

この射は、最大可換代数商の八重有限平坦構造が与える
正規化 trace と、境界 Jacobian 環の次数三部分への Hessian 同型を使う。
中心元の単純な可換化は古典 Casimir の制限と一致しない
（命題\ref{prop:point-casimir-restriction-obstruction}）。
逆構成は次のように強められる。

\begin{maintheorem}[Poisson 対の開部分との同型]
任意の可換 $\C$-代数上の滑らかな三次 $X_H\subset\PP(W)$ について、
各相対双ベクトルは一意な斉次二次 affine 持上げ
$\Pi=\iota_{dH}P$ を持ち、$\iota_{dH}\Pi=0$ を多項式として満たす。
この同定と Poisson 条件の同値は任意基底変更に両立する。

$P$ が非零分解可能な Poisson 対の stack を $\mathfrak P_n^{\mathrm{dec}}$ とする。
周囲代数単独の有限安定化子を要求せず、良い PBW・中心・相対双対・全同値比較の
入力と滑らかな $H_w$ を課した圏 stack $\mathfrak M_n^{\mathrm{sm}}$ には同型
\[
 \mathfrak M_n^{\mathrm{sm}}\simeq
 \mathfrak P_n^{\mathrm{good}}\subset\mathfrak P_n^{\mathrm{dec}}
\]
がある。右辺は非空開であり、両辺は次元 $\binom{n+2}{3}+3n-2$ の
滑らかな既約有限型 DM stack である。逆写像は $\Pi$ の二次関係と
$[w]=\mathfrak d_\Pi^{-1}([H])$ による。表示 gerbe の位数の区分も同じである。

既存の一般三次の被約 Poisson 分類を入力すると、一般三次の非空開上で
これは全非零 Poisson モジュライへの開埋込みとなる。
全 Poisson モジュライのうち滑らかな三次のモジュライを支配する既約成分は一つで、
本構成はその一般部分の開を同定する。
\end{maintheorem}

証明は第\ref{sec:poisson-inverse}節に置く。
これは全 intrinsic 圏モジュライの成分分類でも、Kuznetsov 成分の同値定理でもない。
後者の比較では、むしろ次の障害がある。

\begin{maintheorem}[残余復元との比較障害]
可換復元稿の完備 MF 移送と一般点復元の仮定を満たすとする。
全ての $n\ge3$ に対し、一ブロック中心三次族で
\[
 H^{\mathrm{TH}}=F_0+hxyz,\qquad
 H^{\mathrm{res}}=F_0+
 \frac{h+h^2/3+h^5/12}{1-h^4/4}xyz,\qquad
 F_0=\frac{x^3+y^3+z^3+\sum_{i=1}^{n-1}u_i^3}{3}
\]
となるものがある。二つの滑らかな一般ファイバーは代数閉拡大後も射影同型でない。
従って非可換残余圏と $\operatorname{Ku}(V(H^{\mathrm{TH}}))$ との間には
rotation を保つ dg Morita 同値がない。
$n\not\equiv1\pmod3$ では rotation を指定しない dg Morita 同値もない。
同じ周囲 tensor の滑らかな三次 Casimir の選択だけでは、この障害を解消できない。

一般の混合 Hesse 表示では、短い中心公式から $h^2ABC/6$ を引いて
Hessian 二次 jet を持つ中心元に規格化する。指定した生成元固定の formality 模型の下で
\[
 W_{\mathrm{res}}-H^{\mathrm{TH}}
 =h^2\left(\frac{\mathscr H_F}{24}
       +\frac{\operatorname{tr}(\mathscr H_F)}{12}\right)+O(h^3),
 \qquad \mathscr H_F=\det\operatorname{Hess}_{x,y,z}F.
\]
括弧内の全 Jacobi 環の class は非空 Zariski 開で非零である。
従って特殊ファイバーを恒等に標識した形式比較は、一般の表示係数においても
既に二次 jet で成立しない。
\end{maintheorem}

厳密 Hesse 係数は固定入力の計算であり、次元に一様な比較障害への適用、
固定 tensor の全 Casimir に対する障害、混合係数上の二次 class を本稿で証明する
（第\ref{sec:nfm-residual-th-obstruction}節）。
一般混合係数における有限 jet の障害から、標識を外した全ての一般ファイバーの
圏非同値までを結論してはいない。


\subsection{残余三次に適合する補正}

\begin{maintheorem}[代数的補正と形式的な残余比較]
分解可能な Poisson 対を $(H,P)$ と表し、$K=\operatorname{supp}P$、
$U=\operatorname{ann}K$ とする。$H$ と $H|_{\PP(K)}$ が滑らかな開で、
active 臨界代数の正規化 trace を $\tau_{H,K}$ と置く。
座標不変な三次式
\[
 \mathcal D_P(H)=\langle\det(D_K^2H),P^{\otimes2}\rangle,
 \qquad H^\sharp=H+\frac{\mathcal D_P(H)}{24}
                      +\frac{\tau_{H,K}(\mathcal D_P(H))}{12}
\]
は、非空開で Poisson 対の代数的射 $\Phi:(H,P)\mapsto(H^\sharp,P)$ を定める。
この射は任意底変更と全 O-同値に両立し、始域と終域を適切な非空開へ縮めると、
表現可能な有限 \emph{\'etale} 射で階数 $9$ となる。
補正後の active 平面切断は量子側の楕円点曲線に等しく、
指定した中心正規化では $W_h\equiv H^\sharp\pmod{h^4}$ となる。

一方、滑らかな先頭三次と滑らかな active 制限を持つ一ブロック族について、
固定 formality の生成元固定模型では、形式 Poisson tensor $\pi_h$ は
元の中心一次空間を厳密に保存し、残余 potential $W_h$ に対して
\[
 \pi_h=a_h\,\iota_{dW_h}P_0,\qquad a_h=h+O(h^2)
\]
と一意に書ける。この tensor は周囲二次代数だけから決まる。
その相対化は被約有限型の普遍 PBW 係数模型とその任意底変更で成立する。

二次関係の係数が $h$ に依存せず、Hessian 正規化中心元を用いると、
$\pi_h$ と $a_h$ は奇級数、$W_h$ は偶級数である。
$g=F|_{U=0}$、$S(xyz)=1$ と規格化した四次不変式を $S$ とすると、
全混合係数について
\[
a_h=h+\frac{S(g)}{240}h^5+O(h^7),\qquad
 \pi_h=h\Pi_F+\frac{h^3}{24}\Pi_{\mathscr H_F}
                 +h^5\Pi_{Q_5}+O(h^7)
\]
を得る。また active 制限は
\[
 [h^5]\pi_h|_{U=0}=\frac{S(g)}{80}\Pi_g.
\]
従って $S(g)\ne0$ の稠密開では、固定した $h$ と生成元に関して
代数的 tensor 補正と形式逆 tensor は五次から異なる。
五次 tensor の全混合係数は、明示的な相対三次式 $I(F),J(F)$ を用いて
\[
 \pi_5(F)=\Pi_{Q_5(F)},\qquad
 Q_5(F)=\frac{4I(F)+J(F)-2S(g)F}{240}
\]
と書ける。指定した中心元
$\operatorname{sym}(F-h^2\mathscr H_F/12+h^4\Psi_\lambda(F))$ の残余三次は
\[
 W_h=F+\frac{h^2}{24}\mathscr H_F
       +h^4\left(\Psi_\lambda(F)+\frac{I(F)}{24}-\frac{J(F)}{80}\right)+O(h^6)
\]
である。全 Jacobi 環での有限補正三次との差は
\[
 [h^{-4}(W_h-H^\sharp)]_{h=0}
 =\frac{[I(F)]}{60}+\frac{[\tau_F(I(F))]}{40}
                      +\frac{S(g)[\tau_F(F)]}{80}.
\]
この類は各 $n\ge3$ の混合係数の非空開で非零であり、
$\C[h]/(h^5)$ 上で特殊ファイバーの同定を保つ三次の同型を妨げる。
固定 MF・window 入力の下では、同じ標識を保つ残余圏の同値も妨げる。

一方、式\eqref{eq:nfm-fourth-invariant-correction}の座標不変な三次式
$\mathcal E_P(H)$ による
\[
 (H,P)\longmapsto
 \left(H^\sharp+\mathcal E_P(H),
       \left(1+\frac{S_P(H)}{240}\right)P\right)
\]
は、適切な非空開で任意底変更と全 O-同値へ降下する。
指定中心族に沿い、この出力三次は $W_h$ と $h^6$ を法として、
出力 tensor は $\pi_h$ と $h^7$ を法として一致する。


分離型 $g_\lambda+G(U)$、
$g_\lambda=(x^3+y^3+z^3)/3+\lambda xyz$、$1+\lambda^3\ne0$ では、
$G$ の偏微分の共通零点が原点だけであると仮定する。
量子化尺度 $s$ に対する補正後三次は
\[
 H^{\mathrm{corr}}=g_\lambda+G+
               \frac{s^2}{24}\det\operatorname{Hess}(g_\lambda).
\]
可換復元稿の完備 MF・相対 window・固定標的への有限降下と
coherent 残余圏の飽和性を仮定すると、
全 $n\ge3$ において、ある有限拡大 $K'/\C(s)$ 上で
$(\mathcal K_{B/(w)},\rho)\simeq
(\operatorname{Ku}(V(H^{\mathrm{corr}})),\rho_{\mathrm{corr}})$
が指定三次周期を保って成立する。
\end{maintheorem}

代数的射と次数は定理\ref{thm:nfm-algebraic-hessian-correction}、
\ref{thm:nfm-correction-degree-nine}、形式構成は定理\ref{thm:nfm-residual-formal-poisson}、
jet と五次の相違は系\ref{cor:nfm-mixed-even-jet}、
\ref{cor:nfm-fifth-restriction-obstruction}、混合項の還元は
命題\ref{prop:nfm-mixed-fifth-reduction}、\ref{prop:nfm-mixed-fifth-explicit}、
四次 potential と比較障害は定理\ref{thm:nfm-mixed-fourth-potential}、
系\ref{cor:nfm-corrected-fourth-obstruction}、
追加の代数的射は定理\ref{thm:nfm-fourth-algebraic-correction}、
split 比較は
定理\ref{thm:nfm-split-corrected-residual}に置く。
これらの有限補正式そのものの一般混合係数での全次数比較は主張しない。
内在的な残余三次への全次数比較は、次の定理で別に構成する。
有限 jet の標識付き障害から、一般ファイバーで標識を外した Ku の非同値は結論しない。
従来の $\Phi$ の一般次数九と点曲線の厳密な一致を、
追加補正の射にも自動的に移すことはしない。

\subsection{一般点の代数的残余射と very general な超越性}

\begin{maintheorem}[残余三次と射影 Poisson 対の一般次元での復元]
$n\ge3$ とし、可換復元稿の完備 MF 移送、正則底の window、
一般点生成性と coherent 飽和性を用いる。
良い一ブロック O-marked stack の非空稠密開上で、低 rotation 積
\[
 V_\rho=\operatorname{Nat}^0(\mathrm{id},\rho),\qquad
 Q_\rho=\ker\bigl(\Sym^2V_\rho\to
              \operatorname{Nat}^0(\mathrm{id},\rho^2)\bigr)
\]
から、任意底変更に両立する方程式直線
\[
 \mathcal L=\ker\bigl(\Sym^3V_\rho\to
       V_\rho\otimes(\Sym^2V_\rho/Q_\rho)\bigr)
\]
を復元できる。中心一次作用 $c:U\hookrightarrow V_\rho$ と合わせた
$(\mathcal L,\operatorname{ann}(cU))$ は、尺度を忘れた分解可能 Poisson 対への
全 O-同値に両立する代数的射 $\Psi_{\mathrm{res}}$ を与える。
この射は支配的で、さらに開を縮めれば相対次元一で滑らかである。
有限 \'etale 被覆上で、対応する同じ $X_{\mathrm{res}}=V(\mathcal L)$ に対し
\[
 \mathcal K\simeq\operatorname{Ku}(X_{\mathrm{res}})
\]
が成立し、被覆上で選んだ rotation と指定三次周期も保存できる。

生成元固定の harmonic-angle 模型からの相対 flow を選べば、
一般形式像は $(V(W_h),[\pi_h])$ と同定される。
他方、その実際の $\pi_h$ の固定座標での係数は、
各 $n\ge3$ の very general な滑らかな係数で $\C(h)$ 上代数的でない。
この非代数性は有限代数拡大では解消せず、普遍係数の関数体上でも成立する。
\end{maintheorem}

代数的な族と全同値への降下は定理\ref{thm:nfm-projective-residual-morphism}、
支配性は命題\ref{prop:nfm-projective-residual-dominance}、
比較被覆は定理\ref{thm:nfm-residual-comparison-cover}、
形式像の同定は系\ref{cor:nfm-projective-formal-compatibility}で証明する。
非代数性は定理\ref{thm:nfm-normalized-tensor-transcendence}による。
後者は指定された固定 frame の主張であり、形式座標変更を許した
射影三次の代数化や、任意の intrinsic な尺度の選択を否定しない。
尺度を持つ lift は $\mathcal L\otimes\det K_{\mathrm{res}}$ の非零切断に対応する。
次の定理は、この線束の平方に点曲線から標準切断を与える。

方程式の一意性は比較核の一意性を含まない。
核全体は自己同値 stack の torsor であり、同型類を整合させた後に
スカラー二 cocycle の障害が残る（命題\ref{prop:nfm-residual-kernel-descent}）。
一変数 trait では、形式比較を選んだ滑らかな三次模型の存在も示すが、
その正準性や核の $h=0$ への延長は主張しない
（系\ref{cor:nfm-residual-good-reduction}）。

\subsection{点曲線が定める尺度と一般有限な残余射}

\begin{maintheorem}[楕円点曲線による残余 Poisson 射の規格化]
上と同じ固定入力の下で、良い一ブロック O-marked stack の
非空稠密開 $\mathfrak U$ を取る。
楕円点曲線を $E_{\mathrm{pt}}$、残余 active 曲線を
$C_{\mathrm{res}}=X_{\mathrm{res}}\cap\mathbf P(K_{\mathrm{res}})$ とすると、
全 $n\ge3$ について
\[
 j(E_{\mathrm{pt}})=j(C_{\mathrm{res}}),\qquad
 \mathscr T=\mathcal L\otimes\det K_{\mathrm{res}}
     \simeq\lambda_{\mathrm{res}}^*.
\]
$j=0,1728$ を除けば、点曲線の標準微分は
$\mathscr T^{\otimes2}$ の至る所非零な標準切断を定め、
全 O-同値と任意基底変更に両立する代数的射
\[
 \Psi_{\mathrm{res}}^\pm:\mathfrak U
    \longrightarrow[\mathfrak P_n^{\mathrm{dec}}/\mu_2]
\]
を与える。$\mu_2$ は Poisson tensor の符号を変える。
この射は支配的かつ一般有限で、適切な非空開へ制限すると有限 étale 全射となる。
さらに反対代数の自由な作用 $\iota$ と整な中間被覆 $Z$ により
\[
 \mathfrak U\xrightarrow{\ 2\ }\mathfrak U/\langle\iota\rangle
       \xrightarrow{\ \sim\ }Z\xrightarrow{\ 9\ }Y,
 \qquad \deg\Psi_{\mathrm{res}}^\pm=18
\]
と因子化する。$Y$ は符号商標的の稠密開であり、矢印上は有限 étale 階数を表す。

標準的な有限 étale 階数二の被覆
\[
 \mathfrak D=
 \underline{\operatorname{Isom}}_{\mathrm{gp}}
 (\operatorname{Jac}(E_{\mathrm{pt}}),\operatorname{Jac}(C_{\mathrm{res}}))
 \longrightarrow\mathfrak U
\]
上では尺度を含む射
$\widetilde\Psi_{\mathrm{res}}:\mathfrak D\to\mathfrak P_n^{\mathrm{dec}}$ を得る。
両曲線の平面同型の sheaf $\mathfrak I$ は有限 étale 階数十八であり、
$\mathfrak I\to\mathfrak D$ は $\operatorname{Jac}(E_{\mathrm{pt}})[3]$-torsor である。
点 shift は $\mathfrak D$ 上で残余 active 曲線の平行移動へ自然に移り、
deck 反転はその逆と Poisson tensor の負を同時に与える。
また $L_0=\C(c_4,-c_6)$ を補正前の active 不変量の体とすると、
$\mathfrak D$ の一般点の Kummer 類は、一意な
$\chi_{\mathrm{act}}\in L_0^\times/(L_0^\times)^2$ の引戻しである。
この類は零であり、適切な非空開上で $\mathfrak D$ は分裂する。
二つの切断のうち半古典先頭項で規格化したものは、
全 O-同値に自然な full Poisson 射
\[
 \Psi_{\mathrm{res}}^{\mathrm{sc}}:
 \mathfrak U\longrightarrow\mathfrak P_n^{\mathrm{dec}},
 \qquad\deg\Psi_{\mathrm{res}}^{\mathrm{sc}}=9
\]
を定める。適切な非空開への制限後に有限 \'etale 階数九となり、
$\Psi_{\mathrm{res}}^{\mathrm{sc}}\circ\iota
 =-\Psi_{\mathrm{res}}^{\mathrm{sc}}$ を満たす。
\end{maintheorem}

証明は\S\ref{sec:nfm-residual-point-scale}に置く。
核心となる全次数の恒等式は、生成元固定模型における
\[
 W_h|_{U=0}=b_h\left(g+\frac{h^2}{24}\det\operatorname{Hess}(g)\right),
 \qquad g=F|_{U=0},\quad b_h\in1+hR[[h]]
\]
である。この式は MF 圏の $U=0$ への基底変更を仮定せず、
active 再帰と既存の同座標 Hesse 同定から証明する。
点曲線による tensor は形式 frame で $(h/c_h)\pi_h$ となり、
固定 harmonic tensor と尺度が同じとは主張しない。
また $\mathfrak D$ 上に両 genus-one torsor の同型があるとは限らず、
平面同型にはさらに $\mathfrak I\to\mathfrak D$ を用いる。
正確な次数の証明は、混合係数の全次数三角形と共通の Hesse 安定化群を用いる
（命題\ref{prop:nfm-residual-mixed-triangular}、定理\ref{thm:nfm-residual-exact-degree-eighteen}）。
反対被覆は full Poisson 符号被覆の active な二次 twist であり、その一般点の類は
$\chi_{\mathrm{op}}=\chi_{\mathrm{sgn}}+\chi_{\mathrm{act}}$ を満たす
（系\ref{cor:nfm-opposite-cover-active-twist}）。
Hesse 係数体への二次類の写像は単射である
（補題\ref{lem:nfm-hesse-kummer-injective}）。
混合係数に関する三角多項式は、三捩れ対称点でも滑らかな三次を与える。
標的の有限対角と、指定された一般点同型の正規基底への延長により、
その近傍の Poisson 対を代数的に降下する
（補題\ref{lem:nfm-poisson-target-finite-diagonal}、
\ref{lem:nfm-finite-diagonal-rational-extension}）。
延長上の三捩れ自己同型から有理三捩れ点を得て Jacobian の二次 twist を排除し、
$\chi_{\mathrm{act}}=0$ と半古典規格化された一般次数九の full 射を得る
（定理\ref{thm:nfm-active-sign-vanishing}、系\ref{cor:nfm-semiclassical-full-residual}）。
従って $\chi_{\mathrm{op}}=\chi_{\mathrm{sgn}}$ である。
この延長の特殊ファイバーを低 rotation 積の特殊化から再構成できることや、
比較核がそこへ延びることまでは主張しない。
一般の反対代数の対は O-marked 圏として異なるが、同じ残余三次と、
共通の有限被覆上で同値な残余 dg 圏を持つ
（系\ref{cor:nfm-opposite-same-residual-category}）。

比較核については、split な表示 chart で stabilization の混合 Euler 数を
可換復元の入力における三次 Euler 数に合わせる開かつ閉な数値型を選べる。
その中で核の同型類が整合的に選ばれた後の scalar gerbe の類 $\alpha$ は
\[
 d_n\alpha=0,\qquad
 \operatorname{ind}(\alpha_{\mathrm{gen}})\mid d_n,
 \qquad d_n=3^{\lfloor(n+1)/2\rfloor}
\]
を満たす（命題\ref{prop:nfm-kernel-three-primary}）。
試験対象を置くための表示 gerbe を分裂させる必要がある体上では、
$n\equiv1\pmod3$ のときに限り上界にさらに因子三を加えられる
（系\ref{cor:nfm-generic-kernel-bound}）。
この bound は比較核の同型類を選ぶ前段階を解決せず、
標準二重被覆だけに核が降りることも意味しない。
むしろ既に選んだ同型類の scalar index は、三に素な体拡大では不変である
（系\ref{cor:nfm-kernel-prime-to-three-invariance}）。
全 rotation 軌道の Euler pairing の整数イデアルも正確に $d_n\mathbf Z$ であり、
その反復だけでは bound を改善できない。
選んだ厳密相対 flow は完備中心環の斉次係数商に降り、active MF 比較を与える
（命題\ref{prop:nfm-completed-central-restriction}）。
有限の代数的比較核にこの coherent な中心線型性を降ろす段階は未完成である。

\begin{maintheorem}[active 制限による比較核の三捩れ評価]
残余復元と比較被覆に用いた固定入力を仮定する。
split 表示の一般点 $\Spec K$ 上で、原点剰余体の stabilization $\kappa_{\mathrm{nc}}$ と、
可換側の係数商・window $1$ からなる関手
\[
 q_Y:\operatorname{Ku}(X_{\mathrm{res}})
       \longrightarrow\operatorname{Perf}(C_{\mathrm{res}}),
 \qquad C_{\mathrm{res}}=X_{\mathrm{res}}\cap\PP(K_{\mathrm{res}})
\]
を固定する。一般点の有限分離拡大上に、rotation と指定周期を保つ同値核 $P$ で、
\[
 R_P=q_Y\Phi_P(\kappa_{\mathrm{nc}}),\qquad
 H^{i}(R_P)=0\ (i>-1),\qquad
 H^{-1}(R_P)\text{ が次数零の線束}
\]
を満たすものが存在する。

この型に属する比較核の同型類が $K$ 上で整合的に選ばれ、その実現の
$\mathbf G_m$-gerbe を $\mathcal G$、類を $\alpha$ とする。
すると $\mathcal G$ 上には重み一・階数三の束があり、
\[
 3\alpha=0,\qquad \operatorname{ind}(\alpha)\mid3
\]
である。その endomorphism 代数は $K$ 上の次数三の Azumaya 代数へ降下し、
符号を除いて $\alpha$ を表す。
\end{maintheorem}

証明では、次数一の中心正則列を殺した剰余体の導来像を、内部次数の purity により
非正準に分裂させる。形式比較の後に負の rotation を計算すると、
最高次 cohomology に線束が一つだけ残る。
この cohomology 自体は正準であり、選んだ full kernel の gerbe 上に重み一で降下する。
これを実際の次数三線束 $\OO_{C_{\mathrm{res}}}(1)$ と tensor して押し下げると、階数三束を得る。
型の非空性は命題\ref{prop:nfm-kernel-active-line-nonempty}、
index の証明は定理\ref{thm:nfm-kernel-active-line-bound}に置く。
さらに、この最高次線束の相対 Picard 類を $m_P$ とすると、核の gerbe は
Picard gerbe の $m_P$ による引戻しに同値である
（系\ref{cor:nfm-kernel-picard-obstruction}）。
より強く、scalar 障害は $J_C(S)/3J_C(S)$ を通り、$m_P=3q$ なら消える。
一般にも三分点被覆
\[
 T_P^{(3)}=S\mathbin{\times}_{m_P,J_C,[3]}J_C\longrightarrow S
\]
は階数九の有限 \'etale 全射であり、その上に最高次線束を規格化した
full 比較核が一意な同型を除いて一意に存在する
（定理\ref{thm:nfm-kernel-tripling-cover}、系\ref{cor:nfm-kernel-picard-mod-three}）。
これは選択済みの full 同型類と、固定した表示・window に関する構成である。
また active 曲線上では実際の比較核が存在し、体上では一般直線の切片が
次数三の有限 étale な scalar 分裂被覆を与える
（系\ref{cor:nfm-kernel-active-curve-splitting}）。
表示が非分裂な体では、表示分裂後にこの型となるという仮定の下で、
$n\not\equiv1\pmod3$ なら三、$n\equiv1\pmod3$ なら九が index の上界となる
（系\ref{cor:nfm-kernel-active-line-unsplit}）。
この強い型の評価を、従来の Euler 数だけを指定した型の全ての核へは拡張しない。
また核の同型類の選択や gerbe の中立性は含まない。
active 同値核そのものの選択や coherent な中心線型性の代数的降下は仮定せず、
full kernel から得る一つの対象を用いている。

さらに支配的な完備族上で相対 THCR/window1 比較と底変更を構成する
（命題\ref{prop:nfm-relative-active-thcr}、\ref{prop:nfm-relative-active-window-basechange}）。
従って同じ固定入力の下で、少なくとも一つの支配的な有限 étale 被覆上に、
平面同型 $\varphi:E_{\mathrm{pt}}\simeq C_{\mathrm{res}}$ と比較核 $P$ で
\[
 \begin{gathered}
 \operatorname{Jac}(\varphi)=\theta_{\mathrm{sc}},\qquad
 m_P=\theta_{\mathrm{sc}}(p_{\mathrm{pt}})=p_{\mathrm{res}},\\
 \alpha=\delta_{C_{\mathrm{res}}}(p_{\mathrm{res}})
 \end{gathered}
\]
となる型が存在する。最後の式はこの型の full 同型類を選んだ後の主張である
（命題\ref{prop:nfm-kernel-point-line-nonempty}、
系\ref{cor:nfm-kernel-point-picard-obstruction}）。
$\theta_{\mathrm{sc}}$ は半古典規格化による Jacobian 同型であり、
$\varphi$ はその上の階数九の $J_E[3]$-torsor $\mathfrak I_{\mathrm{sc}}$ から選ぶ。
証明の核は、楕円曲線上の $U_{a,b}=t_a^*\circ(-\otimes P_b)$ に対する
rotation の共役式であり、右 point shift を織り込むと $3a=0$、$b=p_{\mathrm{pt}}$ となる。
同じ部分群の rotation 保存自己同値は Picard 類を変えないため、
これだけで $m_P$ を任意に零へ調整することはできない。
この非空性は、全ての比較核への適用や full 同型類の選択を含まない。

この障害は実際の active 点対応において一般に非零となる。

\begin{maintheorem}[一般点の平面同型と比較核の Picard 障害]
良い split 一ブロック表示の標準一般係数体を $k_n$ とする。
残余復元と完備 MF・window 等の固定入力の下で、任意の $n\ge3$ に対して
\[
 \mathfrak I_{\mathrm{sc}}(k_n)\ne\varnothing.
\]
従って半古典同型 $\theta_{\mathrm{sc}}$ を誘導する偏極曲線同型
$E_{\mathrm{pt}}\simeq C_{\mathrm{res}}$ が元の係数体上に存在する。
生成元固定形式模型の恒等分枝に一致する同型を選ぶと、これは表示変更に共変であり、
良い一ブロック O-marked stack の共通な滑らかな稠密開上の切断
$\varphi_{\mathrm{sc}}\in\mathfrak I_{\mathrm{sc}}(\mathfrak U)$ へ降下する。
この切断は全 O-保存同値と任意の底変更に両立する。

一般 split 係数体の actual active 点曲線とその shift 点について
\[
 \begin{gathered}
 \delta_{C_{\mathrm{res}}}(p_{\mathrm{res}})
   =\delta_{E_{\mathrm{pt}}}(p_{\mathrm{pt}}),\\
 \operatorname{per}\delta_{E_{\mathrm{pt}}}(p_{\mathrm{pt}})
   =\operatorname{ind}\delta_{E_{\mathrm{pt}}}(p_{\mathrm{pt}})=3
 \end{gathered}
\]
となる。$p_{\mathrm{pt}}$ と $p_{\mathrm{res}}$ はそれぞれの Jacobian の
$k_n$-有理点群で三倍ではない。

前述の点型の実際の full 比較核 $P$ が有限拡大 $L/k_n$ 上に存在すれば、
$3\mid[L:k_n]$ である。
従ってこの型の核は元の体および次数が三に素な有限拡大上には存在しない。
この不存在には、元の体での full 同型類の事前選択は不要である。

点型 full 核の同型類を $k_n$ 上で選択済みの場合、その scalar 類は
\[
 \alpha=\delta_{C_{\mathrm{res}}}(p_{\mathrm{res}})
       =\delta_{E_{\mathrm{pt}}}(p_{\mathrm{pt}})
\]
であり、$\operatorname{per}\alpha=\operatorname{ind}\alpha=3$ である。
この scalar 段階には三次分離分裂体があり、最小次数は三である。
\end{maintheorem}

平面同型の分裂は定理\ref{thm:nfm-generic-plane-isom-splitting}で証明する。
まず全次数模型の混合次数が有限であることから、平面同型 torsor の類を
active 係数体 $K=\C(g_\alpha,h)$ へ降ろす。
次に $H_g=\operatorname{Stab}_{\operatorname{SL}_3}(g)$ の全 ambient 上の作用を、
混合項が零の滑らかな対称点まで延長する。
\[
 Z_{\operatorname{PGL}(X\oplus U)}(H_g)\simeq\operatorname{GL}(U)
\]
によりその点の active 平面同型が元の体上に存在し、残った定数 torsor が消える。
最後に全係数を $\widetilde F=hF$ と変数変換すると体は $k_n(h)$ であり、
有限 \'etale scheme の有理点が純超越拡大で増えないことから $k_n$ へ戻せる。
特殊点で比較核を延長することや、形式完備化が三捩れ類を検出することを仮定しない。
表示変更への共変性と stack 上への降下は
命題\ref{prop:nfm-semiclassical-plane-isom-descent}による。
純中心項も含めた TH potential の parabolic 共変性を用い、
せん断による項を落とさずに形式恒等分枝を比較する。
点型の表示非依存性は系\ref{cor:nfm-point-type-O-descent}で別に証明する。

点側の非零性は命題\ref{prop:nfm-active-point-sklyanin-realization}、
\ref{prop:nfm-active-point-obstruction-dominant}と
系\ref{cor:nfm-generic-point-picard-index-three}で示す。
出発点は $\C(a,b,\lambda)$ 上の
$X^3-aY^3-bZ^3+\lambda XYZ=0$ と三捩れ平行移動の障害
$\bigl(a,b+\lambda^3/(27a)\bigr)_3$ の直接計算である。
Jacobian の関数体上で無限位数の点に替えた後、
Sklyanin potential を元の体上へ降下する。
支配性を検査する有限 Hesse 拡大と、非零類を保持する実際の純超越拡大とを区別する。
表示係数体におけるこの障害は、複素数体上の各点で非零という意味ではない。

残余側への移送と full 核への帰結は
系\ref{cor:nfm-generic-residual-picard-index-three}、
\ref{cor:nfm-point-comparison-cover-three-divisibility}による。
この全次元の系は点型を扱う。$n=3,4$ では、以下の自己同値分類を介して
全 full 核へ障害を拡張し、同型類選択を含めて最小次数三を決定する。
一般の二曲線については式\eqref{eq:nfm-picard-plane-transfer}の cup 補正があり、
それは二曲線の genus-one torsor の差だけに依存する
（補題\ref{lem:nfm-picard-torsor-quotient}）。
形式的には恒等に合同な平面同型を持ちながら、この補正が非零となる例もある
（命題\ref{prop:nfm-cyclic-plane-twist}）。
今回の一般点では $[\mathfrak I_{\mathrm{sc}}]=0$ なので、その相殺は起こらない。

active 曲線間の比較については、さらに完全な存在判定がある。
実次数三偏極 $(E,L_E),(C,L_C)$ 上で
$F\rho_{E,p}\simeq\rho_CF$ となる Fourier--Mukai 同値が存在することは、
偏極曲線同型の存在と $\delta_E(p)=0$ に同値である
（系\paperref{main}{cor:nfm-elliptic-any-intertwiner-criterion}）。
数値作用を制限しないこの判定から、一般点でそのような active 同値は存在しない。
同型類の相対 sheaf は $\mathfrak I_\theta$、実現 gerbe は
$\delta_C(\theta(p))$ の引戻しである
（命題\ref{prop:nfm-elliptic-intertwiner-gerbe}）。
full 核から active 核を作る代数的 assignment は未構成なので、
この非存在を full 核そのものの非存在と同一視しない。
この active 分類からの推論と、以下で与える full 核の独立した降下証明を区別する。

\subsection{点型の内在性、低次元の full 降下と三分点被覆}
\begin{maintheorem}[三・四次元の full 比較と周期標識の最小次数]
\label{main:nfm-low-dimensional-full-descent}
固定入力の下で、良い標準 split 一ブロック表示の一般係数体 $k_n$ を考える。
$n=3,4$ なら、幾何的な full 比較核の全ての同型類を $k_n$ 上で選択できる。
各類の実現 gerbe は band を保って互いに同値であり、その共通の Brauer 類は
\[
 \alpha=\delta_{C_{\mathrm{res}}}(p_{\mathrm{res}}),
 \qquad\operatorname{per}(\alpha)=\operatorname{ind}(\alpha)=3
\]
である。従って、点型かどうかにかかわらず実際の full 核は $k_n$ 上にも
三に素な次数の有限拡大上にも存在しない。
一方、ある三次分離体上には実際の点型 full 核が存在する。
これは rotation を intertwine する性質を持ち、
指定周期の追加標識を忘れた full 比較の最小体次数は正確に三である。

両側の三次方程式の生成元から定まる指定周期を加えると、
全 full 同型類に共通な $\mu_3$-torsor が追加される。
source の生成元を保ち、target の生成元を適切な $k_n$-単元倍に取り替えれば、
この torsor を分裂させて標識付き最小次数を三にできる。
一方、target の生成元だけを独立変数 $t$ 倍した $k_n(t)$ 上の一般族では、
標識付き full 比較核の任意の有限定義体の次数は九の倍数となり、最小次数は九である。
\end{maintheorem}

詳細は補題\ref{lem:nfm-projector-kernel-uniqueness}、
\ref{lem:nfm-cyclic-kernel-class-descent}と
定理\ref{thm:nfm-low-dimensional-full-class-index-three}にある。
幾何的自己同値群は実際の $k_n$-線型核 $T=\rho^{-1}[1]$ が生成する $\mathbb Z$ で、
$T^2=\rho$、$T^3=[1]$ となる。
比較核の類の Galois cocycle は連続準同型 $\Gamma_{k_n}\to\mathbb Z$ なので零である。
$T$ との合成が scalar gerbe の band を保つため、全類の障害が一致する。
指定周期に適合する rotation 同型の torsor の Kummer 類を
$\gamma\in k_n^\times/k_n^{\times3}$ とする。
命題\ref{prop:nfm-period-phase-class-invariance}により、$\gamma$ は全 full class に共通である。
source と target の生成元をそれぞれ $u,v$ 倍すると
$\gamma\mapsto[v/u]\gamma$ となる
（命題\ref{prop:nfm-period-equation-frame-change}）。
独立尺度 $t$ の場合、$\gamma=[d]$ と書けばその成分体は
$k_n(t)[s]/(s^3-dt)=k_n(s)$ である。
この純超越体上では $\alpha$ の index 三が保たれるため、
全体の最小次数は九となる（定理\ref{thm:nfm-free-period-frame-degree-nine}）。
両側の同時規格化では位相は変わらない。
元の復元が与える target のデータは方程式直線であるが、
命題\ref{prop:nfm-period-equation-cube-frame}により full tensor と
指定平面同型から生成元を三乗倍まで定められる。
一般係数体ではこの平面同型自体も唯一である
（系\ref{cor:nfm-generic-plane-isom-unique}）。
全 $n\ge3$ の正準強化 class についても、
命題\ref{prop:nfm-adjoint-period-all-dimensions}により
$\gamma_{\mathrm{sc},n}$ が一意に定義され、形式比較体への像は自明である。
混合有理 CF 比較が同じ完備核へ移り、忠実な体拡大から随伴復元条件が反映されるので、
全 class の分類を使わずにこの正準 class を追跡できる。
さらに定理\ref{thm:nfm-period-phase-active-descent}と
命題\ref{prop:nfm-period-phase-invariant-descent}により、これは一意な類
\[
 \overline\gamma_n\in\C(s,t)^\times/\C(s,t)^{\times3},
 \qquad s=c_4(g),\quad t=-c_6(g),\qquad n\ge3
\]
から来る。
有限 MF の内部次数と正の混合重みにより無限 arity の挿入後も混合単項式数を抑え、
$\kappa(m)$ 上に実際の周期付き full Morita 比較を作ることが最初の段階である。
次に線型持上げと周期解を $(f,\eta)\mapsto(\lambda f,\lambda^{-1}\eta)$ で商にし、
周期 torsor 自身の $\operatorname{SL}_3$ 同変構造を得る。
一般 Heisenberg 安定化群の character が消えるため、不変量体への降下が成立する。
これは形式体への Kummer 引戻しが単射だと仮定した結果ではない。

$\Delta=s^3-t^2$、$D=1-s/24-t/216-s^2/6912$ と置く。
Hesse 係数 $q=2+z$ の境界における計算は、次の値を与える。

\begin{maintheorem}[正準周期位相の Hesse 境界値]
\label{main:nfm-period-hesse-boundary}
全 $n\ge3$ の標準 split 一般係数体 $k_n$ 上で、上の固定入力と
三乗まで規格化した方程式を用いる。
$m$ を混合係数の組とし、$R_H=\C(p,m)[[z]]$、
$K_H=\C(p,m)((z))$ を $q=2+z$ に対応する完備 Hesse trait とする。
このとき
\[
 \left.\gamma_{\mathrm{sc},n}\right|_{K_H}=[D_R],\qquad
 D_R=-\frac{(p^3+8)^2}{27}z+O(z^2).
\]
従って $v_D(\overline\gamma_n)=1\pmod3$ であり、
位相被覆はこの境界で分岐指数三を持つ。
全ての正準位相 $\gamma_{\mathrm{sc},n}$ は非自明である。
$\gamma_{\mathrm{sc},n}=[a_n]$ に対し
$M_n=k_n(\sqrt[3]{a_n})$ と置けば $[M_n:k_n]=3$ で、
正準強化 class の指定周期を標識した最小実現次数は
\[
 3\operatorname{ind}(\alpha_{M_n})\in\{3,9\},
 \qquad \alpha=\delta_{C_{\mathrm{res}}}(p_{\mathrm{res}})
\]
となる。$n=3,4$ では、この値は全 full class に共通する。
\end{maintheorem}

局所 Kummer 類の証明は命題\ref{prop:nfm-boundary-period-local-kummer}、
最小次数の公式は系\ref{cor:nfm-adjoint-period-marked-minimum}による。
Hom の平坦性だけでなく、二次関係の飽和、三次積の最初の非零項、
指定した一般 active 表示同型の延長まで確認することで、
付値と剰余単元の双方を比較している。

この torsor が $\Spec\C[s,t,1/(\Delta D)]$ 上で有限 \'etale に延びるという
追加仮定の下では、補題\ref{lem:nfm-period-two-boundary-criterion}の
三候補 $1,[D],[D]^2$ から上の付値が $[D]$ を選ぶ。
従って大域式 $\overline\gamma_n=[D]$ はこの非分岐仮定の下で成立する。
その仮定と異なる次元の大域的な $\overline\gamma_n$ の一致は未証明である。
残る位相の計算は、$\overline\gamma_n[D]^{-1}$ の他の因子での分岐と、
位相体 $M_n$ 上の $\alpha$ の index を対象とする。
後者が正準規格化での標識付き最小次数三／九を決める。

全次元では、原点対象と active 制限に内在的な表示
\[
 \kappa_{\mathrm{nc}}\simeq j_{\mathrm{nc}}^!(\mathcal O)[1],
 \qquad q_Y\simeq\mathbf Li_C^*j_Y
\]
を与える。従って点型は
$Q_P=\mathbf Li_C^*j_Y\Phi_Pj_{\mathrm{nc}}^!(\mathcal O)$ の
最高次層 $H^0(Q_P)$ が $p_{\mathrm{res}}$ を表すという条件となり、
系\ref{cor:nfm-point-type-O-descent}により全 O-同値に適合する。
点型と $\alpha=\delta_C(p_{\mathrm{res}})$ は共通な良い稠密開上で表示非依存である。
最高次線束型の選択済み類には、非分裂 ambient でも
$\alpha=\delta_C(m_P)$ と $3[C]=0$ から $3\alpha=0$ が従うが、
index は既存の三または九の上界と区別する
（系\ref{cor:nfm-unsplit-point-scalar-period-three}）。

\begin{maintheorem}[随伴復元点型の正準 class と最小次数]
\label{main:nfm-finite-adjoint-point-classes}
全 $n\ge3$ の標準 split 一般係数体 $K=k_n$ で、固定入力の仮定を保つ。
$Z=j_{\mathrm{nc}}^!\mathcal O$、$q=\mathbf Li_C^*j_Y$、
$q^R=j_Y^!i_{C,*}$ と置く。
点型核の最高次線束 $M_P=H^0(q\Phi_PZ)$ に対する自然な随伴射
\[
 \Phi_PZ\longrightarrow q^RM_P
\]
が同型となる型は非空である。この型の幾何的 full class は、
有限 \'etale 群 scheme $H_p$ の torsor $S_p$ をなし、その実現 stack は
\[
 S_p\times_K\mathscr P_{p_{\mathrm{res}}}
\]
に同値である。ここで $\mathscr P_{p_{\mathrm{res}}}$ は残余曲線上の
$p_{\mathrm{res}}$ を表す線束の Picard gerbe である。
一般一ブロック族では全 $n\ge3$ で
\[
 H_p=1,\qquad S_p=\operatorname{Spec}K.
\]
従って class は一意である。その scalar gerbe はなお
$\alpha=\delta_{C_{\mathrm{res}}}(p_{\mathrm{res}})$ を持ち、
$\operatorname{per}(\alpha)=\operatorname{ind}(\alpha)=3$ である。
周期同型そのものを標識しないこの型の最小有限実現次数は正確に三となる。
\end{maintheorem}

非空性と有限 torsor の記述は
定理\ref{thm:nfm-adjoint-point-finite-class-torsor}、
最小次数の一般公式は系\ref{cor:nfm-adjoint-point-realization-degree}にある。
全次元の自明性は定理\ref{thm:nfm-origin-rotation-faithful}と
系\ref{cor:nfm-generic-adjoint-class-unique}による。
原点と内部 shift を保つ kernel の作用を全 $A_\infty$ 構造と homotopy まで回復するので、
三連続の生成対象の対象像だけから忠実性を推論してはいない。
元の最高次線束条件だけの全点型や、任意の full 核の分類は主張しない。

\begin{maintheorem}[非分裂点型の指数九と普遍階数の最小性]
\label{main:nfm-unsplit-picard-index-nine}
$n\ge4$、$n\equiv1\pmod3$ とし、固定入力の下で点型と半古典平面同型が
全 $\mathcal O$-降下する共通な良い一ブロック開を取る。
その任意の非空開に、体値点であって
\[
 \operatorname{per}\delta_{C_{\mathrm{res}}}(p_{\mathrm{res}})=3,
 \qquad\operatorname{ind}\delta_{C_{\mathrm{res}}}(p_{\mathrm{res}})=9
\]
となるものが存在する。この点で実際の点型 full 比較核を持つ有限体拡大の次数は、
九の倍数である。表示安定化子が自明な共通稠密開内では、
全該当次元で随伴復元点型の最小実現次数は九である。

普遍 Picard gerbe $\mathscr P_{p_{\mathrm{res}}}$ 上には重み一・階数九の束があり、
底をどの非空開へ縮めても重み一の正の階数は九の倍数である。
表示 gerbe の index 三と、この点の Picard index 九は異なる。
全ての体値点で index 九となること、
共通な良い開の外での class 選択は主張しない。
型を指定しない全比較核については次の定理が別の最小次数を与える。
\end{maintheorem}

証明は命題\ref{prop:nfm-unsplit-genus-one-picard-nine}、
\ref{prop:nfm-unsplit-point-index-nine-family}、
系\ref{cor:nfm-universal-unsplit-picard-rank-nine}、
\ref{cor:nfm-unsplit-point-minimum-nine}にある。
独立な二つの三次 Kummer 座標から、
次数九の division algebra を二重付値で明示する。
この九は周期標識を追加しない点型の scalar 障害から来る。

\begin{maintheorem}[非分裂 rotation と低次元の全型の最小次数]
\label{main:nfm-unsplit-full-minimum}
最高次線束型の full class $P$ を選択し、その障害を
$\alpha=\delta_C(m_P)$、次数三偏極類の障害を $\beta$ とする。
幾何的 rotation の class は Galois 不変で、その実現 gerbe は $\beta$ である。
$\rho^rP$ の障害は $\alpha+r\beta$ であり、$r=\pm1$ なら index は三を割る。

$n=3,4$ で幾何学的自己同値群が $\mathbf Z$ となる範囲では、
$\alpha=\delta_C(p_{\mathrm{res}})$ として全型の最小有限実現次数は
\[
 d_{\mathrm{full}}(K)=\min_{r=0,1,-1}\operatorname{ind}(\alpha+r\beta)
 =\begin{cases}1,&\alpha\in\langle\beta\rangle,\\
 3,&\alpha\notin\langle\beta\rangle.
 \end{cases}
\]
特に前定理の四次元の指数九族では、強化点型の最小次数は九、全型では三である。
周期同型そのものを標識する場合は含まない。
\end{maintheorem}

証明は補題\ref{lem:nfm-unsplit-rotation-gerbe}、
定理\ref{thm:nfm-unsplit-rotation-index-three}、
\ref{thm:nfm-unsplit-lowdim-full-degree}、
系\ref{cor:nfm-unsplit-nine-versus-three}にある。
高次元にも class 選択後の三次上界は有効だが、
rotation で型が変わり得るので、点型自身の指数九は保持される。

五次元では、Clifford 表示の半 rotation は
\[
 O_\Pi\simeq T=\rho^{-1}[1],\qquad O_\Pi^3\simeq[1]
\]
と perfect 核として同定できる
（定理\ref{thm:nfm-fivefold-half-rotation-identification}）。
従って既存の実核降下で得た $O$ も元の体上で $T$ に等しく、
候補 $\epsilon=OT^{-1}$ は恒等である。
この具体的同定と、上記の全次元の安定化群の忠実性は独立の論証である。

\begin{maintheorem}[標準三分点被覆の正確な次数と Galois 群]
\label{main:nfm-third-division-monodromy}
全 $n\ge3$ の同じ一般係数体で
$[3]^{-1}(p_{\mathrm{pt}})\simeq[3]^{-1}(p_{\mathrm{res}})$ は連結な次数九の体である。
一つの三分点を定義する最小次数は九、
Galois 閉包の群は $\mathbf F_3^2\rtimes\operatorname{SL}_2(\mathbf F_3)$、次数は二百十六となる。
この九次体は非 Galois で、真の中間体を持たない。
一つの三分点を持つ最小の Galois 拡大の次数も二百十六である。
\end{maintheorem}

証明は定理\ref{thm:nfm-generic-point-third-division-monodromy}による。
最小三分点次数九は、上の最小 scalar 分裂次数三と異なる。
従って三次分裂体上でも点は三で割れず、九次体自身に三次の中間体はない。
任意の full 核が三分点を選ぶという主張は含まない。
$n\ge5$ でも強化点型の class は一意に選べるが、
全型への非存在障害の拡張は未解決である。

\subsection{Jacobi スキームの次元境界}
\begin{maintheorem}[Jacobi スキームの境界と平方零構造]
滑らかな可換三次 $X_F\subset\PP^{n+1}$ の自然な Jacobi スキームについて、
$3\le n\le7$ の一般三次では
\[
 \operatorname{Pois}(X_F)\simeq\widehat{\operatorname{Gr}}(3,n+2)
\]
である。
$n\ge8$ では任意の滑らかな $F$ についてこの等式は成立せず、
頂点完備化も抽象局所環として同型でない。
Poisson 幾何の研究稿\cite{NCCubicA}の一般三次に対する被約分類を入力すると、この場合の一般の
$\operatorname{Pois}(X_F)$ は非被約であり、
冪零根基の次数二部分の次元は少なくとも
\[
 \epsilon_N=5\binom{N+1}{6}-N\binom N4+\binom N3>0\quad(N=n+2\ge10)
\]
である。$n=8$ ではさらに
\[
 J_{10}/I_F\simeq\C(-2)^{330},\qquad
 H_{R_{10}/I_F}(t)=H_{R_{10}/J_{10}}(t)+330t^2
\]
となる。冪零根基は極大イデアルで消える平方零イデアルであり、
この拡大は完備局所環としても分裂しない。
$n=8$ のこの結論は、既存の被約分類を用いない独立な計算証明による。
さらに全ての $n\ge3$ の一般三次について
\[
 \operatorname{Pois}(X_F)\setminus\{0\}
 =\widehat{\operatorname{Gr}}(3,N)\setminus\{0\},\qquad
 \operatorname{Proj}(R_N/I_F)=\operatorname{Gr}(3,N)
\]
がスキームとして成り立つ。
$n\ge8$ では冪零根基は頂点に支持される非零の有限長加群であり、
頂点の局所環は深さ零である。
一方、全ての $n\ge11$ で、非零 Poisson 点の局所環に
330次元の平方零部分空間を持つ滑らかな特殊三次が存在し、射影化しても非被約である。
\end{maintheorem}

Jacobi スキームの主張は\S\ref{sec:nfold-jacobi-scheme}で証明する。
Jacobi イデアル $I_F$ の $N\binom N4$ 個の生成二次式の間には
$\binom N3$ 個の普遍的な独立従属関係がある。
Pl\"ucker イデアル $J_N$ の二次部分の次元 $5\binom{N+1}{6}$ との比較が上の下界を与える。
$N\le10$ の最大階数には、明示的な整数小行列が法 $32003$ で非零である証明書を付す。
$n=8$ では Pl\"ucker の線形 syzygy を330次元の商へ射影し、
$39600\times39600$ 小行列を法二で検証することにより、次数三以降の差を消す
（定理\ref{thm:nfold-eightfold-square-zero}）。
標数二の三次幾何を使うのではなく、元の整数係数行列の階数を特性零へ移す。
非零分解可能 $P$ の support を $A$ とし、
$K_{F,A}=\{w\in W:T_F(A,w,W)=0\}$ と置くと、法方向の接空間は
$\Lambda^3K_{F,A}$ となる（定理\ref{thm:nfold-normal-tangent}）。
これが一般三次での非零部分の一致と頂点支持を与える。
さらに、support を固定した横断切片そのものを $\Lambda^3K_{F,A}$ 上の
polar Jacobi スキームとして記述する（定理\ref{thm:nfold-transverse-slice}）。
Thom--Sebastiani 直和にはスキームの retract があり、
これを用いて特殊三次の非零点へ冪零構造を移す。
高次数の Hilbert 関数の比較から、冪零部分の最高次数と
頂点の極大イデアルによる消滅指数が次元によらず一様に有界ではないことも示す
（命題\ref{prop:nfold-higher-nilradical}）。
非被約性の結論に用いる被約分類の依存と、方程式数の比較だけで従う結論を分ける。
また、これらの Poisson スキームを全 Abel 圏の Artin 変形空間と同一視したわけではない。

